\documentclass[11pt]{article}
\usepackage{../common/polystyle}
\usepackage{xr}
\externaldocument[P-]{poly_pacbayes}
\graphicspath{{figures/}}
\renewcommand{\polytitle}{PAC-Bayesian Learning -- Solutions to the exercises}
\newcommand{\statement}[1]{\begin{tcolorbox}[enhanced,breakable,colback=black!3,colframe=black!25,boxrule=0.4pt,sharp corners,fontupper=\itshape\small]#1\end{tcolorbox}}
\newcommand{\exo}[1]{\Needspace{6\baselineskip}\subsection*{Exercise #1}\addcontentsline{toc}{subsection}{Exercise #1}}
\newcommand{\sol}{\noindent\textbf{Solution.}\ }

\begin{document}
\sloppy
\begin{center}
{\LARGE\bfseries PAC-Bayesian Learning}\\[0.2cm]
{\Large Detailed solutions to the exercises of the lecture notes}\\[0.3cm]
{\large Master MALIA -- Ensemble Methods}\\[0.4cm]
Guillaume Metzler\\
Institut de Communication (ICOM), Université Lumière Lyon 2\\
Laboratoire ERIC UR 3083, Lyon, France
\end{center}

\bigskip
\noindent This document gathers detailed solutions to the exercises that close each section of the lecture notes
\emph{PAC-Bayesian Learning: A Theoretical Framework for Ensemble Methods}. Exercises, figures, equations and results
are numbered as in the notes. The numerical values were computed with \texttt{pbtools.py}; the short scripts used
for the simulation exercises are those of the notebooks. We advise to try each exercise before reading its
solution: most of them take a few lines once the right tool of the notes has been identified, and the remarks
that follow the solutions point to what the exercise is meant to teach. The later exercises of each section are
labelled \textbf{(Comprehension)} (short conceptual questions), \textbf{(Application)} (computations by hand or with
\texttt{pbtools.py} on concrete numbers) or \textbf{(Proof)} (short proofs of results used in the notes).

\tableofcontents
\bigskip

%% ======================================================================
\section*{Section 1 -- Why PAC-Bayes for ensemble methods?}
\addcontentsline{toc}{section}{Section 1 -- Why PAC-Bayes for ensemble methods?}

\exo{1.1}
\statement{Write the predictors of bagging, of AdaBoost and of a kernel machine with random Fourier features in the
form \eqref{P-eq:mv}. In each case, what is the set of voters $\Hcal$, and what is the distribution $Q$?}

\sol In all three cases we look for a set of voters $\Hcal$ and a probability distribution $Q$ on $\Hcal$ such that the
predictor equals $\sign(\E_{h\sim Q}h(x))$.

\emph{Bagging and random forests.} The voters are the $n$ trees $h_1,\dots,h_n$ learned on the bootstrap samples, so
$\Hcal=\{h_1,\dots,h_n\}$, and $Q$ is the uniform distribution $Q(h_t)=\frac1n$. Indeed
$\sign\big(\frac1n\sum_th_t(x)\big)=\sign(\E_{h\sim Q}h(x))$.

\emph{AdaBoost.} The voters are the weak learners $h_1,\dots,h_T$, and $Q(h_t)=\alpha_t/\sum_s\alpha_s$. The weights are
non-negative because each weak learner has a weighted error $\varepsilon_t<\frac12$, so that
$\alpha_t=\frac12\ln\frac{1-\varepsilon_t}{\varepsilon_t}>0$, and dividing the score $\sum_t\alpha_th_t(x)$ by the positive
constant $\sum_s\alpha_s$ does not change its sign. If some $\alpha_t$ were negative, we would replace $h_t$ by $-h_t$ and
$\alpha_t$ by $|\alpha_t|$: a vote with signed weights is always a vote with non-negative weights over the set
$\Hcal\cup(-\Hcal)$.

\emph{Random Fourier features.} The kernel machine is $\sign\big(\sum_{k=1}^K\alpha_k\cos(\omega_k^\top x+b_k)\big)$. The voters are
the functions $h_{\omega,b,s}(x)=s\cos(\omega^\top x+b)$ with $s\in\{-1,+1\}$ (which take values in $[-1,1]$ rather than
$\{-1,+1\}$; for binary voters take $s\,\sign\cos(\omega^\top x+b)$ as in \cref{P-fig:rff}), and $Q$ puts the weight
$|\alpha_k|/\sum_j|\alpha_j|$ on the voter $(\omega_k,b_k,\sign\alpha_k)$. The set of voters is infinite, indexed by
$(\omega,b,s)$, and the distribution from which the frequencies are drawn (the Fourier transform of the kernel,
times the uniform distribution of $b$ on $[0,2\pi]$) plays the role of a \emph{prior}: this is the viewpoint of
\citet{letarte2019pseudo} discussed in \cref{P-sec:linear}.

\emph{Remark.} The exercise shows that the question ``which predictor?'' reduces, for all these methods, to the
question ``which distribution $Q$ over voters?'', which is precisely the object of PAC-Bayes.

\exo{1.2}
\statement{Show that the empirical Rademacher complexity of $\mathrm{conv}(\Hcal)$ equals the one of $\Hcal$.
\emph{Hint:} for a fixed vector of signs $\sigma$, the function $f\mapsto\frac1m\sum_i\sigma_if(x_i)$ is linear in $f$.}

\sol Recall that $\widehat{\mathfrak R}_S(\mathcal F)=\E_\sigma\big[\sup_{f\in\mathcal F}\frac1m\sum_i\sigma_if(x_i)\big]$, where the
$\sigma_i$ are independent uniform signs. Since $\Hcal\subseteq\mathrm{conv}(\Hcal)$, the supremum over the convex hull is at
least the supremum over $\Hcal$, so $\widehat{\mathfrak R}_S(\Hcal)\leq\widehat{\mathfrak R}_S(\mathrm{conv}(\Hcal))$.

For the converse, fix $\sigma$ and let $L_\sigma(f)=\frac1m\sum_i\sigma_if(x_i)$, which is linear in $f$. Any
$f\in\mathrm{conv}(\Hcal)$ can be written $f=\sum_{j=1}^kQ_jh_j$ with $h_j\in\Hcal$, $Q_j\geq0$ and $\sum_jQ_j=1$. By linearity,
\[
L_\sigma(f)=\sum_jQ_jL_\sigma(h_j)\leq\max_jL_\sigma(h_j)\leq\sup_{h\in\Hcal}L_\sigma(h).
\]
Taking the supremum over $f$ and then the expectation over $\sigma$ gives
$\widehat{\mathfrak R}_S(\mathrm{conv}(\Hcal))\leq\widehat{\mathfrak R}_S(\Hcal)$, hence equality. (For infinite convex combinations
the same argument applies with integrals.)

\emph{Remark.} Averaging voters does not increase the capacity measured by the Rademacher complexity. This is the
key to the margin bounds of \citet{schapire1998boosting} and \citet{koltchinskii2002empirical}, which do not depend on
the number of voters in the vote.

\exo{1.3}
\statement{In \cref{P-fig:adaboost}, the training error reaches zero after a few dozen rounds. Using the right panel,
explain why a bound based on the empirical margin distribution can keep decreasing afterwards.}

\sol A margin bound (\cref{P-sec:margins}) has the form
$\RD(\BQ)\leq\P_S(M_Q\leq\theta)+O\big(\sqrt{\ln|\Hcal|\ln m/(m\theta^2)}\big)$ for every $\theta>0$. Its empirical term is not the
training error, which is $\P_S(M_Q\leq0)$, but the proportion of training examples whose normalized margin is below
$\theta$. Once the training error is zero, AdaBoost keeps increasing the margins of the examples: the right panel of
\cref{P-fig:adaboost} shows that the empirical distribution function of the margins moves to the right between $T=50$ and
$T=400$ rounds. For a fixed $\theta$, the quantity $\P_S(M_Q\leq\theta)$ therefore keeps decreasing; equivalently, we can
use a larger $\theta$ for the same empirical term, which decreases the complexity term in $1/\theta^2$. The bound thus
keeps improving although the training error no longer moves, which is consistent with the stable (and even slightly
decreasing) test error of the middle panel.

\exo{1.4}
\statement{\textbf{(Comprehension)} True or false? Justify in one or two sentences. (a)~Since the vote of the running
example is better than an average tree, the risk of a majority vote is always smaller than the risk of its Gibbs
classifier, $\RD(\BQ)\leq\RD(\GQ)$. (b)~The Gibbs risk $\RD(\GQ)$ always lies between the smallest and the largest risk
of the voters in the support of $Q$.}

\sol (a)~\emph{False.} The vote errs where at least half of the weight errs, and it can be worse than an average voter.
Take five voters with the uniform posterior and a distribution $\D$ such that, on half of the examples, exactly three
voters err ($W_Q=0.6$) and, on the other half, none errs ($W_Q=0$). Then $\RD(\GQ)=\E W_Q=0.3$ whereas the vote errs on
every example of the first half, $\RD(\BQ)=0.5$. What \emph{is} always true is the first-order bound
$\RD(\BQ)\leq2\RD(\GQ)$ (Proposition~\ref{P-prop:fo}): the vote can be worse than its voters, but not by more than a factor~$2$.
Whether it is better depends on how the errors are spread over the examples (\cref{P-chap:mv}).

(b)~\emph{True.} $\RD(\GQ)=\E_{h\sim Q}\RD(h)$ is an average of the risks of the voters, hence lies between their minimum and
their maximum over the support of $Q$ (see Exercise~1.8). The Gibbs classifier is never better than the best voter,
which is why a vote, not a Gibbs classifier, is what we deploy.

\exo{1.5}
\statement{\textbf{(Comprehension)} A student proposes the following trick: ``since the bound \eqref{P-eq:mcall_intro}
holds for any prior, I learn my posterior $Q$ and then take $P=Q$, so that $\KL(Q\|P)=0$ and the certificate is simply
$\RS(\GQ)+\sqrt{\ln(2\sqrt m/\delta)/(2m)}$.'' Explain what is wrong.}

\sol The bound holds for any prior \emph{chosen before seeing the sample} $S$ on which $\RS(\GQ)$ is computed. The order of
the quantifiers matters: ``for every fixed $P$, with probability $1-\delta$ over $S$, for all $Q$''. A prior $P=Q_S$ built
from the data is not fixed: the proof of Theorem~\ref{P-thm:general} needs to exchange $\E_{S'}$ and $\E_{h\sim P}$ with $P$
independent of $S'$, and this step fails. With $P=Q_S$, the trick would certify \emph{any} learned posterior, even a
Dirac on the classifier with the smallest empirical risk among millions, with the width of a fixed-classifier bound:
this is exactly the naive bound of \cref{P-fig:uniformity}, which is violated in almost every run. The term $\KL(Q\|P)$ is
precisely the price of having chosen $Q$ with the data; it cannot be removed by relabelling the posterior as a prior.

\emph{Remark.} The prior may depend on data that are \emph{not} used in the bound (\cref{P-sec:prior}); this is the legitimate
version of the student's idea.

\exo{1.6}
\statement{\textbf{(Application)} AdaBoost produces three weak learners with weighted errors $\varepsilon_1=0.1$,
$\varepsilon_2=0.2$ and $\varepsilon_3=0.3$. Compute their weights $\alpha_t$ and the normalized distribution $Q$. On a point $x$ where
$h_1(x)=+1$ and $h_2(x)=h_3(x)=-1$, what does the vote predict? Is the best weak learner always obeyed?}

\sol With $\alpha_t=\frac12\ln\frac{1-\varepsilon_t}{\varepsilon_t}$,
\[
\alpha_1=\tfrac12\ln9=1.0986,\qquad \alpha_2=\tfrac12\ln4=0.6931,\qquad \alpha_3=\tfrac12\ln\tfrac73=0.4236,
\]
so that $\sum_s\alpha_s=2.2154$ and $Q=(0.496,\ 0.313,\ 0.191)$. On the point $x$,
$\E_{h\sim Q}h(x)=0.496-0.313-0.191=-0.008<0$ (unnormalized score $1.0986-0.6931-0.4236=-0.018$), so the vote predicts $-1$:
the best weak learner is outvoted, narrowly, by the two others. A single voter imposes its prediction whatever the
others say if and only if its weight exceeds one half, $Q(h_t)>\frac12$; here $Q(h_1)=0.496$ falls just short.

\emph{Remark.} The normalization by $\sum_s\alpha_s$ does not change any prediction; it only makes the weights a
probability distribution, the object of PAC-Bayes.

\exo{1.7}
\statement{\textbf{(Application)} Consider the bound \eqref{P-eq:mcall_intro} with $\delta=0.05$ and $m=1000$. Compute its
complexity term for the uniform posterior $Q=P$, and for a Dirac posterior on one of $n=100$ voters with the uniform
prior (for which $\KL(Q\|P)=\ln n$). How many examples are needed for the complexity term of the uniform posterior to be
below $0.02$?}

\sol The confidence term is $\ln\frac{2\sqrt{1000}}{0.05}=\ln1264.9=7.143$. For $Q=P$, $\KL=0$ and the complexity term is
$\sqrt{7.143/2000}=0.0598$. For the Dirac, $\KL=\ln100=4.605$ and the term is $\sqrt{(4.605+7.143)/2000}=\sqrt{11.748/2000}=0.0766$:
selecting one voter among $100$ costs less than two points of certificate. For the last question we need
$\ln(2\sqrt m/0.05)/(2m)\leq0.02^2=4\times10^{-4}$. The left-hand side is decreasing in $m$; solving numerically (bisection
on $m$ with Python) gives $m\geq10\,392$ (the term equals $0.0204$ for $m=10\,000$ and $0.0199$ for $m=10\,500$).

\emph{Remark.} Dividing the complexity term by three requires about ten times more data: this is the slow $1/\sqrt m$
rate of McAllester's bound, slightly worsened by the $\ln\sqrt m$ in the numerator. The kl form of \cref{P-chap:bounds}
does much better when the empirical risk is small.

\exo{1.8}
\statement{\textbf{(Proof)} Let $\Hcal=\{h_1,\dots,h_n\}$ be finite and $Q\in\M(\Hcal)$. Show that
$\RD(\GQ)=\E_{(x,y)\sim\D}[W_Q(x,y)]$ with $W_Q(x,y)=\sum_tQ(h_t)\ind{h_t(x)\neq y}$, and deduce that
$\min_{t:Q(h_t)>0}\RD(h_t)\leq\RD(\GQ)\leq\max_{t:Q(h_t)>0}\RD(h_t)$.}

\sol By definition of the Gibbs risk and linearity of the expectation (a finite sum can always be exchanged with an
expectation),
\begin{align*}
\RD(\GQ)=\sum_{t=1}^nQ(h_t)\RD(h_t)&=\sum_{t=1}^nQ(h_t)\,\E_\D\ind{h_t(x)\neq y}\\
&=\E_\D\Big[\sum_{t=1}^nQ(h_t)\ind{h_t(x)\neq y}\Big]=\E_\D W_Q.
\end{align*}
For an infinite class, the same exchange is Fubini's theorem, which applies because the integrand is non-negative.
For the inequalities, let $T=\{t:Q(h_t)>0\}$, $r_-=\min_{t\in T}\RD(h_t)$ and $r_+=\max_{t\in T}\RD(h_t)$. Since the weights
$Q(h_t)$, $t\in T$, are non-negative and sum to one,
$r_-=\sum_{t\in T}Q(h_t)r_-\leq\sum_{t\in T}Q(h_t)\RD(h_t)\leq\sum_{t\in T}Q(h_t)r_+=r_+$.

\emph{Remark.} The Gibbs risk is linear in $Q$; minimizing it alone would put all the weight on the best voter. This
is the root of the failure of first-order bounds as learning objectives (\cref{P-sec:algo_fo}).

\exo{1.9}
\statement{\textbf{(Proof)} For voters with values in $\{-1,+1\}$, show that $\ind{h(x)\neq y}=\frac{1-y\,h(x)}{2}$, deduce that
the margin $M_Q(x,y)=y\,\E_{h\sim Q}h(x)$ equals $1-2W_Q(x,y)$, and show that if the vote $\BQ$ errs on $(x,y)$, then
$W_Q(x,y)\geq\frac12$ (whatever the convention chosen for $\sign(0)$).}

\sol Since $y,h(x)\in\{-1,+1\}$, the product $y\,h(x)$ equals $1$ if $h(x)=y$ and $-1$ otherwise. In the first case
$\frac{1-yh(x)}2=0=\ind{h(x)\neq y}$, in the second $\frac{1-yh(x)}2=1=\ind{h(x)\neq y}$. Hence $y\,h(x)=1-2\ind{h(x)\neq y}$, and
taking the $Q$-expectation,
\[
M_Q(x,y)=\E_{h\sim Q}\big[y\,h(x)\big]=1-2\,\E_{h\sim Q}\ind{h(x)\neq y}=1-2W_Q(x,y).
\]
Let $s=\E_{h\sim Q}h(x)$, so that $\BQ(x)=\sign(s)$. If $y=+1$ and $\BQ(x)\neq y$, then $\sign(s)=-1$, which requires $s<0$, or
$s=0$ if the convention is $\sign(0)=-1$; in both cases $s\leq0$. If $y=-1$ and $\BQ(x)\neq y$, then $\sign(s)=+1$, hence $s\geq0$.
In all cases $M_Q=ys\leq0$, i.e.\ $W_Q=\frac{1-M_Q}2\geq\frac12$. Therefore $\RD(\BQ)\leq\P_\D(W_Q\geq\frac12)$, with equality up to
the ties $W_Q=\frac12$, which the analysis counts as errors: this is the (conservative) reading of \eqref{P-eq:mv_risk_W}.

%% ======================================================================
\section*{Section 2 -- Mathematical toolbox}
\addcontentsline{toc}{section}{Section 2 -- Mathematical toolbox}

\exo{2.1}
\statement{Compute $\kl(0.1\|0.2)$ and $\kl(0.2\|0.1)$. Is the binary kl symmetric?}

\sol By definition $\kl(q\|p)=q\ln\frac qp+(1-q)\ln\frac{1-q}{1-p}$. Hence
\begin{align*}
\kl(0.1\|0.2)&=0.1\ln\tfrac12+0.9\ln\tfrac98=-0.0693+0.1060=0.0367,\\
\kl(0.2\|0.1)&=0.2\ln2+0.8\ln\tfrac89=0.1386-0.0942=0.0444.
\end{align*}
The two values differ: the binary kl is not symmetric. The interpretation of \cref{P-chap:tools} helps: $\kl(q\|p)$ measures
the surprise of observing a frequency $q$ when the true rate is $p$. Observing $0.1$ errors when the true rate is
$0.2$ is less surprising than observing $0.2$ when the true rate is $0.1$, because a Bernoulli variable of mean $0.2$
fluctuates more than one of mean $0.1$.

\exo{2.2}
\statement{Show that the Gibbs distribution $Q_\phi$ of Lemma~\ref{P-lem:dv} is the only posterior for which the change of
measure is an equality.}

\sol The proof of Lemma~\ref{P-lem:dv} shows that, for every $Q$ with $\KL(Q\|P)<\infty$,
\[
\KL(Q\|P)+\ln\E_Pe^\phi-\E_Q\phi=\KL(Q\|Q_\phi).
\]
The change of measure is therefore an equality iff $\KL(Q\|Q_\phi)=0$. By Gibbs' inequality (the KL is non-negative and
vanishes only for equal distributions, see TD Exercise~4), this happens iff $Q=Q_\phi$. If $\KL(Q\|P)=+\infty$, the
right-hand side of \eqref{P-eq:dv} is infinite while the left-hand side is finite (when $\E_Q\phi$ is), so there is no
equality either. Hence $Q_\phi$ is the unique posterior achieving equality, and the supremum in
\eqref{P-eq:gibbs_variational} is attained only at $Q_\phi$.

\exo{2.3}
\statement{Using the Occam bound in Pinsker form, how many examples are needed to certify a risk below $0.15$ for a
classifier with empirical risk $0.1$ chosen among $10^6$ candidates, with $\delta=0.05$?}

\sol With the uniform prior on $|\Hcal|=10^6$ classifiers, $\ln\frac1{P(h)}=\ln10^6=13.82$, and $\ln\frac1\delta=\ln20=3.00$. The
Pinsker form of Theorem~\ref{P-thm:occam} gives
$\RD(h)\leq0.1+\sqrt{\frac{13.82+3.00}{2m}}$. The certificate is below $0.15$ iff
\[
\sqrt{\frac{16.81}{2m}}\leq0.05\iff m\geq\frac{16.81}{2\times0.0025}=3362.2,
\]
so $m=3363$ examples suffice. \emph{Remarks.} With the kl form, far fewer examples are needed: about $1550$, since the kl
is much tighter for small risks. And the dependence on the number of candidates is only logarithmic: with $10^{12}$
candidates, the numerator becomes $27.63+3.00$ and we need $6127$ examples, less than twice as many.

\exo{2.4}
\statement{Using Pinsker's inequality, show that the kl interval of \cref{P-fig:concentration} is never wider than the
Hoeffding interval. Using notebook~1, determine for which empirical risks it is at least twice narrower when $m=1000$.}

\sol Let $\varepsilon=\ln(1/\delta)/m$. The Hoeffding interval is $[\hat q,\hat q+\sqrt{\varepsilon/2}]$ and the kl interval is
$[\hat q,\klup(\hat q,\varepsilon)]$. If $p$ belongs to the kl interval, then $\kl(\hat q\|p)\leq\varepsilon$, and Pinsker's inequality
gives $2(p-\hat q)^2\leq\kl(\hat q\|p)\leq\varepsilon$, i.e.\ $p\leq\hat q+\sqrt{\varepsilon/2}$. Hence
$\klup(\hat q,\varepsilon)\leq\hat q+\sqrt{\varepsilon/2}$: the kl interval is always contained in the Hoeffding interval.

Numerically, for $m=1000$ and $\delta=0.05$, the Hoeffding width is $\sqrt{\ln20/2000}=0.0387$. Computing
$\klup(\hat q,\varepsilon)-\hat q$ on a grid of $\hat q$, the ratio of the two widths is $3.97$ at $\hat q=0.01$, still $2$ at
$\hat q\approx0.054$, and decreases to $1.00$ at $\hat q=0.5$. The kl interval is thus at least twice narrower for all
empirical risks below about $5.4\%$. Pinsker's inequality is tight around $\frac12$, which is why the two intervals
coincide there.

\exo{2.5}
\statement{\textbf{(Comprehension)} True or false? Justify briefly. (a)~If $P$ and $Q$ are two distributions on a finite
set of voters, $\KL(Q\|P)$ is finite. (b)~If $P$ is uniform on $n$ voters, then $\KL(Q\|P)\leq\ln n$ for every posterior $Q$.
(c)~In the Occam bound with the uniform prior, doubling the number of candidate classifiers doubles the complexity term
$\ln\frac1{P(h)}$.}

\sol (a)~\emph{False.} The KL is finite only if $Q$ is absolutely continuous with respect to $P$, i.e.\ $P(h)=0\Rightarrow Q(h)=0$.
With $P=(1,0)$ and $Q=(\frac12,\frac12)$, the term $Q(h_2)\ln\frac{Q(h_2)}{P(h_2)}$ is infinite: a posterior may not put mass
on a voter that the prior excludes. (It is finite as soon as $P$ has full support: then $\KL(Q\|P)\leq\ln\frac1{\min_hP(h)}$.)

(b)~\emph{True.} $\KL(Q\|P)=\ln n-H(Q)$ and the entropy $H(Q)=-\sum_hQ(h)\ln Q(h)$ is non-negative, since each term
$-Q(h)\ln Q(h)\geq0$. Equality holds exactly for the Dirac posteriors, for which $H=0$.

(c)~\emph{False.} With $N$ candidates, $\ln\frac1{P(h)}=\ln N$; with $2N$ candidates it becomes $\ln N+\ln2$: doubling the class
only adds $\ln2\approx0.69$ nats. This logarithmic growth is visible in \cref{P-fig:occam}.

\exo{2.6}
\statement{\textbf{(Application)} Let $P$ be uniform on $n=4$ voters. Compute $\KL(Q\|P)$ for $Q=(\frac12,\frac14,\frac14,0)$, for
$Q'=(0.7,0.1,0.1,0.1)$ and for the Dirac $\delta_{h_1}$, both from the definition and with the formula $\ln n-H(Q)$.
Compute also $\KL(P\|Q')$ and $\KL(P\|Q)$, and comment.}

\sol \emph{Posterior $Q$.} From the definition (the term with $Q(h_4)=0$ vanishes, with $0\ln0=0$),
$\KL(Q\|P)=\frac12\ln\frac{1/2}{1/4}+2\cdot\frac14\ln\frac{1/4}{1/4}=\frac12\ln2=0.3466$. With the entropy,
$H(Q)=\frac12\ln2+2\cdot\frac14\ln4=1.0397$ and $\ln4-H(Q)=1.3863-1.0397=0.3466$.

\emph{Posterior $Q'$.} $\KL(Q'\|P)=0.7\ln2.8+3\times0.1\ln0.4=0.7207-0.2749=0.4458$; and $H(Q')=0.9404$, so that
$\ln4-H(Q')=0.4458$.

\emph{Dirac.} $\KL(\delta_{h_1}\|P)=1\cdot\ln\frac{1}{1/4}=\ln4=1.3863$, the maximal value, and $H(\delta_{h_1})=0$.

\emph{Reversed divergences.} $\KL(P\|Q')=\frac14\ln\frac{0.25}{0.7}+\frac34\ln\frac{0.25}{0.1}=-0.2574+0.6872=0.4298\neq0.4458$: the KL is not
symmetric. $\KL(P\|Q)=+\infty$, since $P(h_4)=\frac14>0$ while $Q(h_4)=0$.

\emph{Comment.} PAC-Bayes uses $\KL(Q\|P)$, the divergence of the posterior \emph{from} the prior. A posterior may ignore
some voters (here $Q(h_4)=0$) at a finite cost, while the reverse divergence would forbid it. Note also that $Q$, which
discards one voter, is cheaper than $Q'$, which keeps them all but puts $70\%$ of the mass on one: the KL measures
concentration, not the size of the support.

\exo{2.7}
\statement{\textbf{(Application)} Check the two test-set certificates of the running example of this section: a vote with
$4.4\%$ of errors on $540$ test points, certified by Hoeffding's inequality ($0.097$) and by the kl inversion ($0.069$) with
$\delta=0.05$. Recompute both certificates with $\delta=0.01$, and comment on the cost of the extra confidence.}

\sol For a fixed predictor and $n=540$ independent test points, the Hoeffding certificate is
$\hat q+\sqrt{\ln(1/\delta)/(2n)}$ and the kl certificate is $\klup(\hat q,\ln(1/\delta)/n)$ (Proposition~\ref{P-prop:testset}).
With $\hat q=0.044$ and $\delta=0.05$: $\ln20=2.996$, so the Hoeffding width is $\sqrt{2.996/1080}=0.0527$ and the certificate is
$0.0967\approx0.097$; the kl budget is $2.996/540=0.00555$ and $\klup(0.044,0.00555)=0.0690$. (With the exact frequency $24/540=0.0444$
one finds $0.0971$ and $0.0696$.) With $\delta=0.01$: $\ln100=4.605$, the Hoeffding certificate is
$0.044+\sqrt{4.605/1080}=0.044+0.0653=0.1093$ and the kl certificate is $\klup(0.044,4.605/540)=\klup(0.044,0.00853)=0.0761$.

\emph{Comment.} Going from $95\%$ to $99\%$ confidence costs $0.013$ with Hoeffding and only $0.007$ with the kl: the
confidence enters through $\ln(1/\delta)$, so it is cheap, and even cheaper with the kl for a small risk, whose
interval width scales like $\varepsilon$ rather than $\sqrt\varepsilon$ (Lemma~\ref{P-lem:pinsker}).

\exo{2.8}
\statement{\textbf{(Proof)} Prove Markov's inequality (Lemma~\ref{P-lem:markov}) in detail, then deduce Chebyshev's inequality: for
a random variable $Z$ with mean $\mu$ and variance $v$, $\P(|Z-\mu|\geq a)\leq\frac{v}{a^2}$ for every $a>0$. Compare it with the
one-sided Cantelli inequality (Lemma~\ref{P-lem:cantelli}) when $a^2=v$.}

\sol \emph{Markov.} Let $Z\geq0$ and $a>0$. Pointwise, $a\ind{Z\geq a}\leq Z$: if $Z\geq a$ the left-hand side is $a\leq Z$, otherwise
it is $0\leq Z$. Taking expectations (monotonicity), $a\P(Z\geq a)\leq\E Z$. For the second form, if $\E Z=0$ then $Z=0$ almost
surely and there is nothing to prove; otherwise take $a=\E Z/\delta$: $\P(Z>\E Z/\delta)\leq\P(Z\geq\E Z/\delta)\leq\delta$, so
$Z\leq\E Z/\delta$ with probability at least $1-\delta$.

\emph{Chebyshev.} The variable $(Z-\mu)^2$ is non-negative with expectation $v$, and $\{|Z-\mu|\geq a\}=\{(Z-\mu)^2\geq a^2\}$.
Markov's inequality gives $\P(|Z-\mu|\geq a)\leq\frac{\E(Z-\mu)^2}{a^2}=\frac v{a^2}$.

\emph{Comparison.} For $a^2=v$ (a deviation of one standard deviation), Chebyshev gives the trivial bound $1$, while
Cantelli gives $\frac{v}{v+v}=\frac12$ for each tail. More generally $\frac v{v+a^2}<\frac v{a^2}$: for one-sided deviations Cantelli
is always better. The trick of its proof, shifting the variable by a free parameter $t$ before applying Markov, is
the one that turns the tandem bound into the C-bound (Theorem~\ref{P-thm:cctnd}).

\exo{2.9}
\statement{\textbf{(Proof)} Admitting Hoeffding's lemma $\E\,e^{\lambda(p-\hat p)}\leq e^{\lambda^2/(8m)}$ (Lemma~\ref{P-lem:hoeffding}), prove the tail
bound $\P(p-\hat p\geq\varepsilon)\leq e^{-2m\varepsilon^2}$ by optimizing over $\lambda>0$, and deduce that $p\leq\hat p+\sqrt{\ln(1/\delta)/(2m)}$ with
probability at least $1-\delta$.}

\sol Fix $\lambda>0$. Since $u\mapsto e^{\lambda u}$ is increasing, $\{p-\hat p\geq\varepsilon\}=\{e^{\lambda(p-\hat p)}\geq e^{\lambda\varepsilon}\}$, and Markov's
inequality applied to the non-negative variable $e^{\lambda(p-\hat p)}$ gives
\[
\P(p-\hat p\geq\varepsilon)\leq e^{-\lambda\varepsilon}\,\E\,e^{\lambda(p-\hat p)}\leq\exp\Big(-\lambda\varepsilon+\frac{\lambda^2}{8m}\Big).
\]
This holds for every $\lambda>0$, so we may choose the best one (it is a deterministic parameter of the inequality, not
chosen with the data). The exponent $g(\lambda)=-\lambda\varepsilon+\frac{\lambda^2}{8m}$ is a convex parabola, minimal where
$g'(\lambda)=-\varepsilon+\frac{\lambda}{4m}=0$, i.e.\ $\lambda=4m\varepsilon$, where $g(4m\varepsilon)=-4m\varepsilon^2+2m\varepsilon^2=-2m\varepsilon^2$. Hence
$\P(p-\hat p\geq\varepsilon)\leq e^{-2m\varepsilon^2}$. To invert it, choose $\varepsilon$ such that $e^{-2m\varepsilon^2}=\delta$, i.e.\
$\varepsilon=\sqrt{\ln(1/\delta)/(2m)}$: the event $\{p>\hat p+\varepsilon\}$ is included in $\{p-\hat p\geq\varepsilon\}$ and has probability at most $\delta$.

\emph{Remark.} This ``Chernoff method'' (exponentiate, apply Markov, optimize the free parameter) is the template of
all the proofs of these notes; in Theorem~\ref{P-thm:general}, the change of measure is inserted between the first two steps.

\exo{2.10}
\statement{\textbf{(Proof)} Let $Q_1,P_1$ be distributions on a finite set $\Hcal_1$ and $Q_2,P_2$ distributions on a finite set
$\Hcal_2$, with $P_1,P_2$ positive. Show that $\KL(Q_1\otimes Q_2\|P_1\otimes P_2)=\KL(Q_1\|P_1)+\KL(Q_2\|P_2)$, and deduce that
$\KL(Q^2\|P^2)=2\KL(Q\|P)$, the identity used for the tandem bound. Show also that $\KL(\delta_h\|P)=\ln\frac1{P(h)}$ for a Dirac
posterior.}

\sol By definition of the product distributions, $(Q_1\otimes Q_2)(h_1,h_2)=Q_1(h_1)Q_2(h_2)$, and similarly for $P_1\otimes P_2$. The
logarithm of a product is a sum, so (the pairs with $Q_1(h_1)Q_2(h_2)=0$ contribute $0$)
\begin{align*}
\KL(Q_1\otimes Q_2\|P_1\otimes P_2)
&=\sum_{h_1,h_2}Q_1(h_1)Q_2(h_2)\Big[\ln\frac{Q_1(h_1)}{P_1(h_1)}+\ln\frac{Q_2(h_2)}{P_2(h_2)}\Big]\\
&=\sum_{h_1}Q_1(h_1)\ln\frac{Q_1(h_1)}{P_1(h_1)}\underbrace{\sum_{h_2}Q_2(h_2)}_{=1}\\
&\qquad+\sum_{h_2}Q_2(h_2)\ln\frac{Q_2(h_2)}{P_2(h_2)}\underbrace{\sum_{h_1}Q_1(h_1)}_{=1}\\
&=\KL(Q_1\|P_1)+\KL(Q_2\|P_2).
\end{align*}
With $Q_1=Q_2=Q$ and $P_1=P_2=P$, $\KL(Q^2\|P^2)=2\KL(Q\|P)$: a pair of independent voters costs twice the price of a
single voter, which is why second-order bounds have a doubled KL term. Finally, the Dirac $\delta_h$ gives weight $1$ to $h$
and $0$ elsewhere, so $\KL(\delta_h\|P)=1\cdot\ln\frac{1}{P(h)}$: the complexity of the Occam bound is the KL of a Dirac posterior.

%% ======================================================================
\section*{Section 3 -- PAC-Bayesian generalization bounds}
\addcontentsline{toc}{section}{Section 3 -- PAC-Bayesian generalization bounds}

\exo{3.1}
\statement{Starting from Theorem~\ref{P-thm:general} with $\Delta(q,p)=2(q-p)^2$, prove McAllester's bound directly,
without going through the kl.}

\sol The function $\Delta(q,p)=2(q-p)^2$ is jointly convex (it is a convex function of the linear form $q-p$), so
Theorem~\ref{P-thm:general} applies. It remains to bound $\mathcal I_\Delta(m)=\E_{S'}\E_{h\sim P}e^{2m(\widehat R_{S'}(h)-\RD(h))^2}$.
For a fixed $h$, by Pinsker's inequality, $2(\hat p-p)^2\leq\kl(\hat p\|p)$, so
$\E_{S'}e^{2m(\hat p-p)^2}\leq\E_{S'}e^{m\kl(\hat p\|p)}\leq2\sqrt m$ by Maurer's lemma (Lemma~\ref{P-lem:chernoff}, for $m\geq8$).
Hence $\mathcal I_\Delta(m)\leq2\sqrt m$ and, with probability at least $1-\delta$, for all $Q$,
\[
2\big(\RS(\GQ)-\RD(\GQ)\big)^2\leq\frac{\KL(Q\|P)+\ln\frac{2\sqrt m}{\delta}}{m},
\]
hence
\[
\RD(\GQ)\leq\RS(\GQ)+\sqrt{\frac{\KL(Q\|P)+\ln\frac{2\sqrt m}{\delta}}{2m}}.
\]
\emph{Remark.} One can also bound $\E e^{2m(\hat p-p)^2}$ without the kl, using the sub-Gaussian tail
$\P(|\hat p-p|\geq t)\leq2e^{-2mt^2}$ and the formula $\E e^{Z}=1+\int_0^\infty e^z\P(Z>z)\,dz$; this gives a bound of order
$m$ instead of $2\sqrt m$, i.e.\ a slightly worse constant ($\ln(2m)$ instead of $\ln(2\sqrt m)$). This is the form of the
original bounds of McAllester.

\exo{3.2}
\statement{For $m=500$, $\delta=0.05$, $\RS(\GQ)=0.05$ and $\KL(Q\|P)\in\{0,5,20\}$, compute McAllester's and Seeger's bounds,
and compare with \cref{P-tab:bound_anatomy}.}

\sol The confidence term is $\ln\frac{2\sqrt{500}}{0.05}=\ln894.4=6.80$. The complexity per example
$\varepsilon=\frac{\KL+6.80}{500}$ and the two bounds are:
\begin{center}\small
\begin{tabular}{cccc}
\toprule
$\KL(Q\|P)$ & $\varepsilon$ & McAllester $0.05+\sqrt{\varepsilon/2}$ & Seeger $\klup(0.05,\varepsilon)$ \\
\midrule
$0$ & $0.0136$ & $0.132$ & $0.094$ \\
$5$ & $0.0236$ & $0.159$ & $0.112$ \\
$20$ & $0.0536$ & $0.214$ & $0.154$ \\
\bottomrule
\end{tabular}
\end{center}
Compared with \cref{P-tab:bound_anatomy} ($m=1000$, $\RS=0.12$), the gap between the two bounds is larger here: the
complexity term of Seeger's bound is about $30\%$ smaller than McAllester's when $\RS=0.12$, but about $40\%$ smaller when
$\RS=0.05$. This is the fast rate of the kl at work: the smaller the empirical risk, the more the square root of
Pinsker's inequality is wasteful. Halving $m$ roughly doubles $\varepsilon$, which costs Seeger's bound less than
McAllester's for the same reason.

\exo{3.3}
\statement{Using the linear form of Catoni's bound given after Corollary~\ref{P-cor:catoni}, find the value of $C$ that
minimizes it when $\RS(\GQ)=0.12$, $\KL(Q\|P)=5$, $m=1000$ and $\delta=0.05$, and compare the result with
\cref{P-tab:bound_anatomy}.}

\sol Let $A=\frac{\KL+\ln(1/\delta)}m=\frac{5+3.00}{1000}=0.0080$ and $f(C)=\frac{C\RS+A}{1-e^{-C}}$. Setting
$f'(C)=0$ gives $\RS(1-e^{-C})=(C\RS+A)e^{-C}$, i.e.\ $\RS(e^C-1-C)=A$. With $\RS=0.12$, $e^C-1-C=0.0667$; since
$e^C-1-C\approx\frac{C^2}2+\frac{C^3}6$, a first approximation is $C\approx\sqrt{2A/\RS}=0.365$, and solving numerically gives
$C^\star=0.344$ and $f(C^\star)=0.169$. The exact (non-linear) Catoni bound is minimized at $C=0.371$, where it equals
$0.165$. In \cref{P-tab:bound_anatomy}, Catoni's bound optimized on a grid of $30$ values of $C$ (with a union bound, hence
$\ln30$ added to the numerator) gives $0.166$ for $\KL=5$ and Seeger's bound $0.168$: the linear relaxation loses
only $0.004$ with respect to the exact bound, and the grid costs very little. \emph{Warning:} as for $\lambda$ in the
Hoeffding bound, $C^\star$ depends on $\RS$ and $\KL$, hence on the data; using it directly requires a grid and a union bound.

\exo{3.4}
\statement{Reproduce the experiment of \cref{P-fig:uniformity} for $m\in\{100,1000,10000\}$. Show that the failure frequency
of the naive bound does not go to zero when $m$ grows, and explain why by comparing $\sqrt{\ln(1/\delta)/(2m)}$ with
the typical deviation of the minimum of $1000$ empirical risks.}

\sol We draw $1000$ classifiers of true risk $0.5$ (their empirical risks are independent $\mathrm{Bin}(m,\frac12)/m$),
keep the smallest empirical risk $\hat r_{\min}$, and check whether the naive bound $\hat r_{\min}+\sqrt{\ln(1/\delta)/(2m)}$ is
below $0.5$. Over $2000$ repetitions:
\begin{center}\small
\begin{tabular}{cccccc}
\toprule
$m$ & mean of $0.5-\hat r_{\min}$ & naive width & ratio & naive bound wrong & Occam wrong \\
\midrule
$100$ & $0.160$ & $0.122$ & $1.31$ & $99.8\%$ & $0.3\%$ \\
$1000$ & $0.051$ & $0.039$ & $1.32$ & $99.9\%$ & $0.5\%$ \\
$10000$ & $0.016$ & $0.012$ & $1.32$ & $100\%$ & $0.6\%$ \\
\bottomrule
\end{tabular}
\end{center}
The failure frequency does not decrease. The reason is that both quantities scale as $1/\sqrt m$. Each empirical risk
is approximately Gaussian with standard deviation $\frac{1}{2\sqrt m}$, and the minimum of $n$ independent Gaussian
variables lies about $\sqrt{2\ln n}$ standard deviations below their mean. The selection bias is therefore of order
$\frac{\sqrt{2\ln n}}{2\sqrt m}=\frac{1.86}{\sqrt m}$ for $n=1000$ (in practice $\approx1.6/\sqrt m$ at these sample sizes), while
the naive width is $\frac{\sqrt{\ln(1/\delta)/2}}{\sqrt m}=\frac{1.22}{\sqrt m}$. Their ratio does not depend on $m$: the naive bound
is always too narrow by the same factor. The Occam bound has width $\frac{\sqrt{(\ln n+\ln(1/\delta))/2}}{\sqrt m}=\frac{2.22}{\sqrt m}$,
which exceeds the selection bias: the price $\ln n$ is exactly what is needed to absorb the selection. More data do not
cure selection bias; only accounting for the selection does.

\exo{3.5}
\statement{\textbf{(Comprehension)} In the proof of Theorem~\ref{P-thm:general}, which step uses the joint convexity of $\Delta$,
which step uses the independence of the prior from the sample, and which step is the only probabilistic one? Why does the
bound hold for all posteriors although Markov's inequality is applied only once?}

\sol The joint convexity of $\Delta$ is used in the \emph{first} step, Jensen's inequality
$\Delta(\E_Q\RS(h),\E_Q\RD(h))\leq\E_Q\Delta(\RS(h),\RD(h))$, which moves from the Gibbs risks to the risks of individual voters.
The \emph{third} step is the only probabilistic one: Markov's inequality on the random variable
$Z(S)=\E_{h\sim P}e^{m\Delta(\RS(h),\RD(h))}$, followed by the exchange $\E_{S'}\E_{h\sim P}=\E_{h\sim P}\E_{S'}$, which requires $P$ not to
depend on the sample. The first two steps (Jensen and the change of measure) are deterministic inequalities, true for
every sample and every posterior.

The event on which the bound holds, $\{Z(S)\leq\E Z/\delta\}$, does not involve $Q$ at all. On this single event, of probability
at least $1-\delta$, the two deterministic steps can be applied to every $Q$ at once. This is why one application of
Markov's inequality suffices, and why the posterior may be chosen afterwards, in any way.

\exo{3.6}
\statement{\textbf{(Comprehension)} True or false? Justify briefly. (a)~For the same $\RS(\GQ)$, $\KL(Q\|P)$, $m$ and $\delta$, Seeger's
bound is never larger than McAllester's bound. (b)~Doubling the sample size $m$ halves the complexity term of McAllester's
bound. (c)~In Catoni's bound, the parameter $C$ may be chosen after seeing the data, since the bound holds for all
posteriors. (d)~Dividing $\delta$ by $10$ adds about $2.3$ nats to the numerator of the complexity term.}

\sol (a)~\emph{True.} Both use the same budget $\varepsilon=\frac1m[\KL+\ln\frac{2\sqrt m}\delta]$, and Pinsker's inequality gives
$\klup(q,\varepsilon)\leq q+\sqrt{\varepsilon/2}$ (Lemma~\ref{P-lem:pinsker}); McAllester's bound is a relaxation of Seeger's.

(b)~\emph{False.} The term is a square root: doubling $m$ divides it by about $\sqrt2$, and even slightly less because
$\ln(2\sqrt m/\delta)$ grows. For $\KL=0$ and $\delta=0.05$, it goes from $0.0598$ ($m=1000$) to $0.0433$ ($m=2000$), a ratio of $0.72$.

(c)~\emph{False.} $C$ is a parameter of the comparison function $\Delta$, so each value of $C$ corresponds to a different
probabilistic event. The bound holds for all $Q$ with $C$ fixed, not for all $C$; choosing $C$ with the data requires
a grid and a union bound (as in \cref{P-tab:bound_anatomy}).

(d)~\emph{True.} $\ln\frac{2\sqrt m}{\delta/10}=\ln\frac{2\sqrt m}\delta+\ln10$, and $\ln10\approx2.30$.

\exo{3.7}
\statement{\textbf{(Application)} An ensemble of $n=50$ voters has a uniform prior. With $m=500$ examples and $\delta=0.05$, the
posterior $Q$ uniform over the $5$ best voters has an empirical Gibbs risk of $0.08$. Compute $\KL(Q\|P)$ and the
complexity per example $\varepsilon$, then McAllester's bound, the PAC-Bayes-$\lambda$ bound at $\lambda^\star$, Seeger's bound, and
Catoni's bound for $C\in\{0.5,1,2\}$. Which value of $C$ is legitimate to report?}

\sol A posterior uniform on $5$ of $50$ voters has $\KL(Q\|P)=\ln\frac{50}5=\ln10=2.303$. The confidence term is
$\ln\frac{2\sqrt{500}}{0.05}=\ln894.4=6.796$, so $\varepsilon=\frac{2.303+6.796}{500}=0.01820$. Then:
\begin{itemize}
\item McAllester: $0.08+\sqrt{0.01820/2}=0.08+0.0954=0.1754$;
\item PAC-Bayes-$\lambda$: $\lambda^\star=2/(\sqrt{1+2\times0.08/0.01820}+1)=2/(\sqrt{9.792}+1)=0.484$, and the bound is
$\frac{0.08}{1-0.242}+\frac{0.01820}{0.484\times0.758}=0.1551$ (it coincides with the ``quadratic'' form, see Exercise~3.10);
\item Seeger: $\klup(0.08,0.01820)=0.1418$ (\texttt{bound\_seeger(0.08, np.log(10), 500, 0.05)});
\item Catoni, with $A=\frac{\KL+\ln(1/\delta)}m=\frac{2.303+2.996}{500}=0.01060$ and the value $\frac{1-e^{-C\cdot0.08-A}}{1-e^{-C}}$: $0.1254$ for
$C=0.5$, $0.1370$ for $C=1$ and $0.1814$ for $C=2$.
\end{itemize}
Only a value of $C$ fixed \emph{before} seeing the data can be reported with confidence $0.95$. If we compute the three
values and keep the best one, we must use $\delta/3$ for each (Proposition~\ref{P-prop:union}); this gives $0.1307$, $0.1402$ and $0.1835$,
and $0.1307$ can then be reported legitimately. It is still below Seeger's bound: when the KL is small, Catoni's bound
benefits from not paying the $\ln2\sqrt m$ term, as observed in \cref{P-tab:bound_anatomy}. (Numerically, the best
$C$ for this example is $C\approx0.5$.)

\exo{3.8}
\statement{\textbf{(Application)} A posterior has zero empirical Gibbs risk on $m=200$ examples, with $\KL(Q\|P)=3$ and $\delta=0.05$.
Compute Seeger's bound (using its closed form $1-e^{-\varepsilon}$ when $\RS(\GQ)=0$) and McAllester's bound. How many examples does
each of them need to certify a Gibbs risk below $0.05$ (with the same KL)?}

\sol When $q=0$, $\kl(0\|p)=\ln\frac1{1-p}$, so $\kl(0\|p)\leq\varepsilon\iff p\leq1-e^{-\varepsilon}$. Here $\ln\frac{2\sqrt{200}}{0.05}=\ln565.7=6.338$
and $\varepsilon=\frac{3+6.338}{200}=0.04669$. Seeger's bound is $1-e^{-0.04669}=0.0456$, while McAllester's bound is
$\sqrt{0.04669/2}=0.1528$, more than three times larger.

For the second question we solve numerically (the two bounds are decreasing in $m$). Seeger's bound is below $0.05$
from $m=182$ on ($0.0498$ for $m=182$, $0.0500$ for $m=181$). McAllester's bound requires
$\sqrt{(3+\ln(2\sqrt m/0.05))/(2m)}\leq0.05$, which holds from $m=2103$ on. The kl bound needs almost $12$ times fewer examples.

\emph{Remark.} This is the fast rate of \cref{P-fig:bounds} (middle panel): with zero empirical risk, the kl bound
decreases like $1/m$ (indeed $1-e^{-\varepsilon}\leq\varepsilon$), McAllester's like $1/\sqrt m$.

\exo{3.9}
\statement{\textbf{(Proof)} Show that the binary kl is \emph{jointly} convex in $(q,p)\in[0,1]\times(0,1)$, as required by
Theorem~\ref{P-thm:general}. \emph{Hint:} for $a>0$, $b>0$, write $a\ln\frac ab=a\,f\big(\frac ba\big)$ with the convex function $f(u)=-\ln u$,
and show that the perspective $(a,b)\mapsto a\,f(b/a)$ of a convex function is jointly convex.}

\sol \emph{Perspective.} Let $f$ be convex on $(0,\infty)$ and $g(a,b)=a\,f(b/a)$ for $a,b>0$. Take $(a_1,b_1)$, $(a_2,b_2)$, $t\in[0,1]$,
and set $a=ta_1+(1-t)a_2$, $b=tb_1+(1-t)b_2$. Then
\[
\frac ba=\frac{ta_1}{a}\cdot\frac{b_1}{a_1}+\frac{(1-t)a_2}{a}\cdot\frac{b_2}{a_2},
\]
a convex combination of $b_1/a_1$ and $b_2/a_2$ (the two weights are non-negative and sum to one). By convexity of $f$,
$f(b/a)\leq\frac{ta_1}af(b_1/a_1)+\frac{(1-t)a_2}af(b_2/a_2)$, and multiplying by $a>0$ gives $g(a,b)\leq t\,g(a_1,b_1)+(1-t)\,g(a_2,b_2)$.

\emph{Application to the kl.} With $f(u)=-\ln u$, $g(a,b)=-a\ln\frac ba=a\ln\frac ab$, so that
$\kl(q\|p)=g(q,p)+g(1-q,1-p)$. The first term is jointly convex by the above; the second is the composition of $g$ with the
affine map $(q,p)\mapsto(1-q,1-p)$, hence also jointly convex; the sum of two convex functions is convex. At the boundary
$q\in\{0,1\}$, $g$ is extended by $g(0,b)=0=\lim_{a\to0}a\ln\frac ab$, and the convexity inequality passes to the limit.

\emph{Remark.} The same argument shows that $\KL(Q\|P)$ is jointly convex in $(Q,P)$; joint convexity is what allows
Jensen's inequality to move from the Gibbs classifier to the individual voters.

\exo{3.10}
\statement{\textbf{(Proof)} Show that, for fixed $\hat q>0$ and $\varepsilon>0$, the function
$\lambda\mapsto\frac{\hat q}{1-\lambda/2}+\frac{\varepsilon}{\lambda(1-\lambda/2)}$ is minimized on $(0,2)$ at $\lambda^\star=2/\big(\sqrt{1+2\hat q/\varepsilon}+1\big)$, which is
\eqref{P-eq:lambda_star}, and that its minimal value is $\big(\sqrt{\hat q+\varepsilon/2}+\sqrt{\varepsilon/2}\big)^2$.}

\sol Write $f(\lambda)=\frac{u(\lambda)}{v(\lambda)}$ with $u=\hat q\lambda+\varepsilon$ and $v=\lambda(1-\frac\lambda2)=\lambda-\frac{\lambda^2}2>0$ on $(0,2)$. Then
$f'=\frac{u'v-uv'}{v^2}$ with
\[
u'v-uv'=\hat q\Big(\lambda-\frac{\lambda^2}2\Big)-(\hat q\lambda+\varepsilon)(1-\lambda)=\frac{\hat q}{2}\lambda^2+\varepsilon\lambda-\varepsilon=:N(\lambda).
\]
$N$ is increasing on $(0,2)$ ($N'=\hat q\lambda+\varepsilon>0$), with $N(0)=-\varepsilon<0$ and $N(2)=2\hat q+\varepsilon>0$: $f$ decreases then
increases, and its minimum is at the unique root of $N$ in $(0,2)$,
\[
\lambda^\star=\frac{-\varepsilon+\sqrt{\varepsilon^2+2\hat q\varepsilon}}{\hat q}=\frac{\varepsilon\big(\sqrt{1+2\hat q/\varepsilon}-1\big)}{\hat q}
=\frac{\varepsilon\cdot2\hat q/\varepsilon}{\hat q\big(\sqrt{1+2\hat q/\varepsilon}+1\big)}=\frac{2}{\sqrt{1+2\hat q/\varepsilon}+1},
\]
where we multiplied and divided by $\sqrt{1+2\hat q/\varepsilon}+1$. With $\varepsilon=\frac1m[\KL+\ln\frac{2\sqrt m}\delta]$, $2\hat q/\varepsilon$ is the ratio
of \eqref{P-eq:lambda_star}. For the value, let $a=\sqrt{\varepsilon/2}$ and $b=\sqrt{\hat q+a^2}$, so that $1+2\hat q/\varepsilon=b^2/a^2$ and
$\lambda^\star=\frac{2a}{a+b}$, $1-\frac{\lambda^\star}2=\frac b{a+b}$. Then, using $\varepsilon=2a^2$ and $\hat q+a^2=b^2$,
\[
f(\lambda^\star)=\frac{\hat q\lambda^\star+\varepsilon}{\lambda^\star(1-\lambda^\star/2)}
=\frac{2a\big(\hat q+a(a+b)\big)/(a+b)}{2ab/(a+b)^2}=\frac{(a+b)(b^2+ab)}{b}=(a+b)^2 .
\]
This is the ``quadratic'' relaxation of \cref{P-sec:sb_optim}. \emph{Numerical check:} for $\hat q=0.12$ and $\varepsilon=0.0094$
(example of \cref{P-chap:bounds}), $\lambda^\star=0.325$ and $f(\lambda^\star)=(a+b)^2=0.1778$, the value $0.178$ of the notes.

%% ======================================================================
\section*{Section 4 -- Posteriors, priors and the computation of the bounds}
\addcontentsline{toc}{section}{Section 4 -- Posteriors, priors and the computation of the bounds}

\exo{4.1}
\statement{For a finite class and a uniform prior, show that $\KL(Q_\lambda\|P)\to\ln\frac{|\Hcal|}{k}$ when $\lambda\to\infty$, where $k$
is the number of empirical risk minimizers.}

\sol Let $n=|\Hcal|$, $r^\star=\min_h\RS(h)$ and $\mathcal M=\{h:\RS(h)=r^\star\}$, with $|\mathcal M|=k$. With the uniform prior,
\[
Q_\lambda(h)=\frac{e^{-\lambda m\RS(h)}}{\sum_{h'}e^{-\lambda m\RS(h')}}=\frac{e^{-\lambda m(\RS(h)-r^\star)}}{k+\sum_{h'\notin\mathcal M}e^{-\lambda m(\RS(h')-r^\star)}}.
\]
For $h'\notin\mathcal M$, $\RS(h')-r^\star>0$, so the exponentials tend to $0$ when $\lambda\to\infty$. Hence $Q_\lambda(h)\to\frac1k$ for
$h\in\mathcal M$ and $Q_\lambda(h)\to0$ otherwise: $Q_\lambda$ converges to the uniform distribution $U_{\mathcal M}$ on the
minimizers. The KL is a continuous function of $Q$ on the simplex (with $0\ln0=0$) when $P$ has full support, so
$\KL(Q_\lambda\|P)\to\KL(U_{\mathcal M}\|P)=\sum_{h\in\mathcal M}\frac1k\ln\frac{1/k}{1/n}=\ln\frac nk$ (TD Exercise~4). With a unique minimizer we
recover the Occam price $\ln n$; when several classifiers tie, the Gibbs posterior spreads over them and pays less.

\exo{4.2}
\statement{Explain why using the same sample to learn the centre of a Gaussian prior and to compute the empirical term
of the bound is not allowed. What happens if we do it anyway?}

\sol The only probabilistic step of the proof of Theorem~\ref{P-thm:general} is the third one, where we write
$\E_{S'}\E_{h\sim P}e^{m\Delta(\widehat R_{S'}(h),\RD(h))}=\E_{h\sim P}\E_{S'}e^{m\Delta(\cdots)}$ and bound the inner expectation for each
\emph{fixed} $h$. This requires $P$ not to depend on $S'$: if $P=P_S$ is centred on a classifier learned on $S$, the
voters drawn from $P_S$ are correlated with the sample, their empirical risks are biased downwards, and the inner
expectation can no longer be controlled by Maurer's lemma. The theorem simply does not apply.

If we do it anyway, the ``certificate'' can be badly wrong. Take the extreme case where the posterior equals the prior,
both centred on the learned classifier: the KL is zero, and the bound becomes a concentration inequality for a
classifier that was selected with the data---exactly the naive bound of \cref{P-fig:uniformity}, which fails almost surely
when the selection is strong. With a flexible model (for instance a neural network that interpolates the data),
the empirical risk is close to zero, the KL is zero, and the certificate would claim a near-zero risk regardless of
the true one. The correct procedure is to split $S$: learn the prior on $S_1$ and compute the bound on $S_2$
(Proposition~\ref{P-prop:ddprior}, \cref{P-fig:ddprior}).

\exo{4.3}
\statement{In the Monte Carlo certificate, how should $\delta$ and $\delta'$ be split to minimize the final bound when $N=100$?
when $N=10^4$?}

\sol The final certificate holds with probability $1-\delta_{\text{tot}}$ with $\delta_{\text{tot}}=\delta+\delta'$; we fix
$\delta_{\text{tot}}=0.05$ and choose the fraction $\alpha=\delta'/\delta_{\text{tot}}$ given to the Monte Carlo step. Increasing $\alpha$
makes the Monte Carlo inversion $\klup(\hat r_N,\ln(2/\delta')/N)$ tighter but the PAC-Bayesian bound looser (it uses
$\delta=(1-\alpha)\delta_{\text{tot}}$). The two terms enter only through logarithms, so the optimum is flat. On the example of
\cref{P-fig:montecarlo} (\texttt{breast cancer}, $\KL=32$, $m=398$), a grid search over $\alpha$ gives the following. With
$N=100$ samples, the Monte Carlo term is the main source of looseness, since $\ln(2/\delta')/N$ is large; it deserves most
of the budget, and the optimum is $\alpha\approx0.86$ ($\delta'=0.043$, $\delta=0.007$), with a certificate of $0.380$ against $0.386$
for the even split. With $N=10^4$ samples, the Monte Carlo term is small and the PAC-Bayesian term dominates; the
optimum moves to $\alpha\approx0.3$, and the certificate ($0.2591$) is practically the same as with the even split ($0.2594$).
In practice the even split is a sensible default; the gain from optimizing is small, and the split must anyway be
chosen before computing the bound (or on a grid with a union bound). It is more effective to increase $N$.

\exo{4.4}
\statement{\textbf{(Comprehension)} (a)~What is the Gibbs posterior $Q_\lambda$ when all the voters have the same empirical risk?
(b)~With the uniform prior, voter $h$ makes $3$ more training errors than voter $h'$; what is the ratio $Q_\lambda(h)/Q_\lambda(h')$?
(c)~True or false: a prior $P\propto e^{-\beta\RD(h)}$, which depends on $\D$, is forbidden.}

\sol (a)~If $\RS(h)=r$ for all $h$, then $Q_\lambda(h)=\frac{P(h)e^{-\lambda mr}}{e^{-\lambda mr}}=P(h)$ for every $\lambda$: the posterior stays equal
to the prior, its KL is zero, and nothing is learned. The data can only move the posterior away from the prior if they
discriminate between voters.

(b)~$m\RS(h)$ is the number of training errors of $h$, so
$\frac{Q_\lambda(h)}{Q_\lambda(h')}=e^{-\lambda m(\RS(h)-\RS(h'))}=e^{-3\lambda}$ (the uniform prior cancels). For $\lambda=1$, $h$ is $e^3\approx20$ times less
likely than $h'$.

(c)~\emph{False.} The rule is that the prior must not depend on the \emph{sample} used in the bound. A prior defined from the
unknown distribution $\D$ is a fixed (if unknown) object, so the proof of Theorem~\ref{P-thm:general} applies; the difficulty is
only that the bound then contains quantities depending on $\D$, which must be controlled in turn \citep{catoni2007pac}.

\exo{4.5}
\statement{\textbf{(Application)} Four voters have empirical risks $0.10$, $0.12$, $0.15$ and $0.30$ on $m=100$ examples, and the prior is
uniform. For $\lambda\in\{0.1,0.5,2\}$, compute the Gibbs posterior $Q_\lambda$, its empirical Gibbs risk, $\KL(Q_\lambda\|P)$ and Seeger's bound
with $\delta=0.05$. For $\lambda=0.5$, check numerically that the objective \eqref{P-eq:objective_linear} at $Q_\lambda$ equals
$-\frac1{\lambda m}\ln\E_{h\sim P}e^{-\lambda m\RS(h)}$. Compare with the bounds of the uniform posterior and of the Dirac on the best voter.}

\sol The voters make $10$, $12$, $15$ and $30$ errors, and $Q_\lambda(h_t)\propto e^{-\lambda\times\text{errors}_t}$. For $\lambda=0.5$, the weights are
proportional to $e^{-5},e^{-6},e^{-7.5},e^{-15}$, i.e.\ (dividing by $e^{-5}$) to $1$, $0.3679$, $0.0821$ and $4.5\times10^{-5}$, whose sum is
$1.4500$; hence $Q_{0.5}=(0.6897,\ 0.2537,\ 0.0566,\ 0.00003)$. The other temperatures are computed in the same way
(\texttt{softmax(-lam*m*R)} in \texttt{pbtools}):
\begin{center}\small
\begin{tabular}{ccccc}
\toprule
$\lambda$ & $Q_\lambda$ & $\RS(G_{Q_\lambda})$ & $\KL(Q_\lambda\|P)$ & Seeger \\
\midrule
$0.1$ & $(0.391,\,0.320,\,0.237,\,0.053)$ & $0.1288$ & $0.158$ & $0.274$ \\
$0.5$ & $(0.690,\,0.254,\,0.057,\,0.000)$ & $0.1079$ & $0.619$ & $0.253$ \\
$2$ & $(0.982,\,0.018,\,0.000,\,0.000)$ & $0.1004$ & $1.296$ & $0.251$ \\
\midrule
uniform & $(0.25,\,0.25,\,0.25,\,0.25)$ & $0.1675$ & $0$ & $0.320$ \\
Dirac on $h_1$ & $(1,0,0,0)$ & $0.1000$ & $\ln4=1.386$ & $0.252$ \\
\bottomrule
\end{tabular}
\end{center}
\emph{Check for $\lambda=0.5$.} The objective is $\RS+\frac{\KL}{\lambda m}=0.10791+\frac{0.61918}{50}=0.12029$. On the other side,
$\E_Pe^{-\lambda m\RS}=\frac14e^{-5}\times1.4500$, so $-\frac1{50}\ln\big(\frac14e^{-5}\times1.4500\big)=\frac{1.3863+5-0.3716}{50}=0.12029$: the two values
coincide, as stated by Proposition~\ref{P-prop:gibbs_post}.

\emph{Comment.} All the Gibbs posteriors improve on the uniform posterior, which pays for the bad voter $h_4$. With only
$m=100$ examples, the confidence term $\ln\frac{2\sqrt{100}}{0.05}=\ln400=5.99$ dominates the KL (at most $\ln4=1.39$), so the
certificate hardly depends on how concentrated the posterior is: the best temperature is essentially the Dirac, as in the
example of \cref{P-chap:bounds}. Temperature matters more when the number of voters is large.

\exo{4.6}
\statement{\textbf{(Application)} A practitioner has $500$ examples and $100$ voters. Option~A uses the $500$ examples for the bound with
a data-independent uniform prior; the best posterior found has an empirical Gibbs risk of $0.15$ and $\KL(Q\|P)=2$. Option~B
learns a prior on $150$ examples and computes the bound on the remaining $350$; the best posterior has an empirical Gibbs
risk of $0.10$ and $\KL(Q\|P)=1$. Compare Seeger's bounds ($\delta=0.05$). Can the practitioner compute both and report the
smaller one?}

\sol Option~A: $\klup\big(0.15,\frac{2+\ln(2\sqrt{500}/0.05)}{500}\big)=\klup(0.15,0.0176)=0.2248$. Option~B:
$\klup\big(0.10,\frac{1+\ln(2\sqrt{350}/0.05)}{350}\big)=0.1739$ (\texttt{bound\_seeger}). Option~B is better: the informed prior lowers
both the empirical risk and the KL, which more than compensates for the smaller sample used by the bound.

Each option, decided in advance, gives a valid certificate (Option~B by Proposition~\ref{P-prop:ddprior}). But computing both and
reporting the smaller is a data-dependent choice between two bounds, so it requires a union bound
(Proposition~\ref{P-prop:union}): each option must be computed with $\delta/2=0.025$, which gives $0.2279$ for~A and $0.1777$ for~B.
Reporting $0.1777$ is then legitimate; the price of the choice ($\ln2$ nats) is only $0.004$.

\exo{4.7}
\statement{\textbf{(Application)} In the setting of the Gaussian posterior example ($m=398$, $\KL(Q\|P)=32$), suppose that $N=1000$
classifiers sampled from $Q$ have an average empirical risk of $0.08$. With $\delta=\delta'=0.025$, compute the Monte Carlo upper
bound on $\RS(\GQ)$ and the final certificate. Compare with the certificate that would be obtained if $\RS(\GQ)=0.08$
were known exactly.}

\sol The Monte Carlo step uses the budget $\frac{\ln(2/\delta')}{N}=\frac{\ln80}{1000}=0.00438$, and
$\klup(0.08,0.00438)=0.1078$: with probability at least $1-\delta'$ over the draws, $\RS(\GQ)\leq0.1078$. The PAC-Bayesian step
uses $\ln\frac{2\sqrt{398}}{0.025}=7.375$, hence $\varepsilon=\frac{32+7.375}{398}=0.0989$, and the final certificate is
$\klup(0.1078,0.0989)=0.292$, valid with probability at least $1-\delta-\delta'=0.95$. If $\RS(\GQ)=0.08$ were known exactly, the
whole budget could go to the PAC-Bayesian step, and $\klup\big(0.08,\frac{32+\ln(2\sqrt{398}/0.05)}{398}\big)=0.250$ (or $0.252$ with $\delta=0.025$).

\emph{Comment.} The Monte Carlo estimation costs about $0.04$ with $N=1000$, in line with \cref{P-fig:montecarlo}: the
certificate lies between the values reported for $N=100$ ($0.375$) and $N=10^4$ ($0.26$). The large KL ($32$ nats for
$398$ examples) is the main source of looseness here, not the sampling.

\exo{4.8}
\statement{\textbf{(Proof)} Prove the formula $\KL\big(\mathcal N(\mu_1,\sigma^2 I_d)\,\|\,\mathcal N(\mu_0,\sigma^2 I_d)\big)=\frac{\|\mu_1-\mu_0\|^2}{2\sigma^2}$, and
deduce $\KL(Q_{\mathbf w}\|P)=\frac12\|\mathbf w\|^2$ for $Q_{\mathbf w}=\mathcal N(\mathbf w,I_d)$ and $P=\mathcal N(\mathbf 0,I_d)$.}

\sol The densities are $q(w)=(2\pi\sigma^2)^{-d/2}e^{-\|w-\mu_1\|^2/(2\sigma^2)}$ and $p(w)=(2\pi\sigma^2)^{-d/2}e^{-\|w-\mu_0\|^2/(2\sigma^2)}$; the
normalizing constants are equal and cancel in the ratio:
\[
\ln\frac{q(w)}{p(w)}=\frac{\|w-\mu_0\|^2-\|w-\mu_1\|^2}{2\sigma^2}.
\]
Under $Q$, write $w=\mu_1+\sigma\xi$ with $\xi\sim\mathcal N(0,I_d)$. Then $\|w-\mu_1\|^2=\sigma^2\|\xi\|^2$ and
$\|w-\mu_0\|^2=\|\mu_1-\mu_0\|^2+2\sigma\langle\mu_1-\mu_0,\xi\rangle+\sigma^2\|\xi\|^2$. The terms in $\|\xi\|^2$ cancel and
$\E\langle\mu_1-\mu_0,\xi\rangle=0$, so
\[
\KL(Q\|P)=\E_{w\sim Q}\ln\frac{q(w)}{p(w)}=\frac{\|\mu_1-\mu_0\|^2}{2\sigma^2}.
\]
With $\sigma=1$, $\mu_1=\mathbf w$ and $\mu_0=\mathbf 0$, $\KL(Q_{\mathbf w}\|P)=\frac12\|\mathbf w\|^2$: the KL is the $\ell_2$ regularizer of the SVM and of
PBGD (\cref{P-sec:linear}). With different variances, the same computation gives the formula of the stochastic networks
of \cref{P-chap:algos}.

\exo{4.9}
\statement{\textbf{(Proof)} Let $P$ be a prior on a finite set of parameters $\theta$ and $\ell(\theta,(x,y))=-\ln p(y\,|\,x,\theta)$. Using
Proposition~\ref{P-prop:gibbs_post} with $\lambda=1$, show that the Bayesian posterior minimizes $\E_{\theta\sim Q}\sum_i\ell(\theta,(x_i,y_i))+\KL(Q\|P)$ over
$Q$, and that the minimal value is minus the log marginal likelihood $-\ln\E_{\theta\sim P}\prod_ip(y_i\,|\,x_i,\theta)$.}

\sol Let $L(\theta)=\sum_{i=1}^m\ell(\theta,(x_i,y_i))=-\ln\prod_ip(y_i|x_i,\theta)$ (we assume the likelihood positive for at least one $\theta$
with $P(\theta)>0$). Proposition~\ref{P-prop:gibbs_post} is stated for losses in $[0,1]$, but its proof only uses the variational
formula \eqref{P-eq:gibbs_variational}, which holds for any $\phi$ with $\E_Pe^\phi<\infty$, automatically true on a finite set. With
$\phi=-L$,
\begin{align*}
\min_Q\Big\{\E_Q L+\KL(Q\|P)\Big\}&=-\sup_Q\Big\{\E_Q[-L]-\KL(Q\|P)\Big\}\\
&=-\ln\E_{\theta\sim P}e^{-L(\theta)}=-\ln\E_{\theta\sim P}\prod_{i=1}^mp(y_i|x_i,\theta),
\end{align*}
and the minimum is attained at $Q(\theta)\propto P(\theta)e^{-L(\theta)}=P(\theta)\prod_ip(y_i|x_i,\theta)$, which is Bayes' rule. The
objective equals $m$ times \eqref{P-eq:objective_linear} with $\lambda=1$ and the log-loss.

\emph{Remark.} This is the observation of \citet{germain2016pac}: minimizing a PAC-Bayesian objective over $Q$ yields the
Bayesian posterior, and the optimal value is the negative log evidence, so comparing models by their marginal
likelihood amounts to comparing their optimal PAC-Bayesian objectives.

%% ======================================================================
\section*{Section 5 -- From the Gibbs classifier to the majority vote}
\addcontentsline{toc}{section}{Section 5 -- From the Gibbs classifier to the majority vote}

\exo{5.1}
\statement{Give an example of three voters and a distribution on three points for which $\RD(\BQ)=0$ while $\RD(\GQ)=\frac13$.
Compute the three oracle bounds.}

\sol Let $\D$ be uniform on three points $z_1,z_2,z_3$, and let voter $h_i$ err on $z_i$ only; $Q$ is uniform on
$\{h_1,h_2,h_3\}$. On each point exactly one voter out of three errs, so $W_Q(z_j)=\frac13$ for all $j$. The vote is correct
everywhere ($W_Q<\frac12$): $\RD(\BQ)=0$, while $\RD(\GQ)=\E W_Q=\frac13$. The second-order quantities are
$e_Q=\E W_Q^2=\frac19$ and $d_Q=\E[2W_Q(1-W_Q)]=2\cdot\frac13\cdot\frac23=\frac49$. Hence:
\[
2\RD(\GQ)=\tfrac23,\qquad 4e_Q=\tfrac49,\qquad
\mathcal C_Q=1-\frac{(1-2/3)^2}{1-8/9}=1-\frac{1/9}{1/9}=0
\]
for the first-order bound, the tandem bound and the C-bound respectively.
The C-bound is exact here: $W_Q$ is constant, its variance is zero, and Cantelli's inequality says that a variable with
zero variance equal to $\frac13$ is never above $\frac12$. The example is a small Condorcet jury with perfectly
``anti-correlated'' errors: the disagreement ($\frac49$) exceeds the Gibbs risk ($\frac13$), so Proposition~\ref{P-prop:compare}
predicts the order C-bound $\leq$ tandem $\leq$ first order, as observed.

\exo{5.2}
\statement{Show that the C-bound is always smaller than $1$ when $\RD(\GQ)<\frac12$, and that it equals $4\RD(\GQ)(1-\RD(\GQ))$ when
$\dD(Q)=0$. Compare with the first-order bound.}

\sol Write $R=\RD(\GQ)$ and $d=\dD(Q)$. By Proposition~\ref{P-prop:decomp}, $0\leq d\leq2R(1-R)<\frac12$, so $1-2d>0$. When $R<\frac12$,
$(1-2R)^2>0$, hence $\frac{(1-2R)^2}{1-2d}>0$ and $\mathcal C_Q<1$. (In fact $1-2d\geq1-4R+4R^2=(1-2R)^2$, so the ratio is at most $1$
and $\mathcal C_Q\geq0$.) When $d=0$, $\mathcal C_Q=1-(1-2R)^2=4R-4R^2=4R(1-R)$. Compared with the first-order bound $2R$:
$4R(1-R)\geq2R\iff2(1-R)\geq1\iff R\leq\frac12$. So when the voters never disagree, the C-bound is \emph{worse} than the
first-order bound, by a factor $2(1-R)$ which is close to $2$ for small $R$. This is consistent with
Proposition~\ref{P-prop:compare} ($d=0<R$): without diversity, second-order information can only hurt, since all voters
predict the same thing and the vote is just one of them.

\exo{5.3}
\statement{$(\star)$ Prove that the minimum over $\mu$ of the Chebyshev--Cantelli family equals the C-bound (TD, Exercise~8).}

\sol Let $R=\E W_Q<\frac12$, $e=\E W_Q^2$, $v=e-R^2=\Var W_Q$, and $f(\mu)=\frac{e-2\mu R+\mu^2}{(1/2-\mu)^2}$ for $\mu<\frac12$. The
numerator equals $\E(W_Q-\mu)^2=v+(R-\mu)^2$. Differentiating,
\[
f'(\mu)=\frac{2(\mu-R)(\tfrac12-\mu)+2\big(v+(R-\mu)^2\big)}{(1/2-\mu)^3}=\frac{2\big[v-(R-\mu)(\tfrac12-R)\big]}{(1/2-\mu)^3},
\]
using $(\mu-R)(\frac12-\mu)+(R-\mu)^2=(R-\mu)(R-\frac12)$. The bracket is increasing in $\mu$ and vanishes at
$\mu^\star=R-\frac{v}{1/2-R}=\frac{R/2-e}{1/2-R}$; $f$ decreases before $\mu^\star$ and increases after, so $\mu^\star$ is the minimum
(and $\mu^\star<R<\frac12$ is admissible). At $\mu^\star$, $R-\mu^\star=\frac{v}{a}$ with $a=\frac12-R$, and
$\frac12-\mu^\star=a+\frac va=\frac{a^2+v}{a}$. Hence
\[
f(\mu^\star)=\frac{v+v^2/a^2}{(a^2+v)^2/a^2}=\frac{v(a^2+v)}{(a^2+v)^2}=\frac{v}{v+a^2}.
\]
Finally, with the margin $M_Q=1-2W_Q$: $v=\frac14\Var M_Q$ and $a^2=\frac14(\E M_Q)^2$, so
$f(\mu^\star)=\frac{\Var M_Q}{\Var M_Q+(\E M_Q)^2}=1-\frac{(\E M_Q)^2}{\E M_Q^2}=\mathcal C_Q$, the C-bound \eqref{P-eq:cbound}. The tandem
bound is $f(0)$, hence always at least the C-bound, which proves the first claim of Proposition~\ref{P-prop:compare}.

\exo{5.4}
\statement{In Condorcet's model with $n$ independent voters of error $p$, compute $\eD(Q)$ and $\dD(Q)$, and show that the
C-bound decreases like $1/n$ while the tandem bound tends to $4p^2$.}

\sol Since $nW_Q\sim\mathrm{Bin}(n,p)$, $\E W_Q=p$ and $\Var W_Q=\frac{p(1-p)}n$. Hence
$\eD(Q)=\E W_Q^2=p^2+\frac{p(1-p)}n$ and, by Proposition~\ref{P-prop:decomp}, $\dD(Q)=2(p-\eD(Q))=2p(1-p)\big(1-\frac1n\big)$. The tandem
bound is $4\eD(Q)=4p^2+\frac{4p(1-p)}n\to4p^2$: it never goes to zero, even though the risk of the vote does (for $p=0.3$
it tends to $0.36$). For the C-bound,
\begin{align*}
1-2\dD(Q)&=1-4p(1-p)+\frac{4p(1-p)}n=(1-2p)^2+\frac{4p(1-p)}n,\\
\mathcal C_Q&=1-\frac{(1-2p)^2}{1-2\dD(Q)}=\frac{4p(1-p)/n}{(1-2p)^2+4p(1-p)/n}.
\end{align*}
When $n\to\infty$, $\mathcal C_Q\sim\frac{4p(1-p)}{(1-2p)^2}\cdot\frac1n$. For $p=0.3$ the equivalent is $\frac{0.84}{0.16\,n}=\frac{5.25}n$; the exact values are $0.174$ for $n=25$ and $0.049$ for
$n=101$ (\cref{P-fig:condorcet}). The true risk decreases exponentially fast (by Hoeffding's inequality,
$\RD(\BQ)\leq e^{-2n(1/2-p)^2}$), so even the C-bound is pessimistic, but it captures the qualitative effect of independence.
The tandem bound misses it because it applies Markov's inequality to $W_Q^2$, i.e.\ it measures the deviations of
$W_Q$ from $0$: since $W_Q$ concentrates around $p>0$, $\E W_Q^2\to p^2$ does not vanish. The C-bound measures the deviations
from the well-chosen centre $\mu^\star$ (Exercise~5.3), and only pays for the variance $\frac{p(1-p)}n$, which vanishes.

\exo{5.5}
\statement{\textbf{(Comprehension)} True or false? Justify briefly. (a)~If $\RD(\GQ)<\frac12$, then $\RD(\BQ)<\frac12$. (b)~If $\dD(Q)=0$, all the
voters in the support of $Q$ make the same predictions on almost every $x$. (c)~Estimating the joint error $\eD(Q)$ requires
labels, but estimating the disagreement $\dD(Q)$ does not. (d)~For a fixed Gibbs risk, the tandem bound decreases when the
disagreement increases.}

\sol (a)~\emph{False.} Suppose that $W_Q=0.55$ on $60\%$ of the examples and $W_Q=0$ elsewhere (e.g.\ $20$ voters with the uniform
posterior, $11$ of which err on the same $60\%$ of the examples). Then $\RD(\GQ)=0.33<\frac12$ but $\RD(\BQ)=0.6$. The only general
guarantee is $\RD(\BQ)\leq2\RD(\GQ)$, which is below $1$ but not below $\frac12$.

(b)~\emph{True.} $\dD(Q)=\E_{(h,h')\sim Q^2}\P_x(h(x)\neq h'(x))$ is an average of non-negative numbers; it vanishes iff
$\P_x(h(x)\neq h'(x))=0$ for $Q^2$-almost every pair (for a finite ensemble, for every pair of voters with positive weight). The
vote then coincides almost surely with each of its voters, and $\RD(\BQ)=\RD(\GQ)$.

(c)~\emph{True.} The joint error is the probability that two voters \emph{both err}, which requires $y$; the disagreement only
compares $h(x)$ and $h'(x)$, so it can be computed on unlabeled examples (\cref{P-chap:advances}).

(d)~\emph{True.} By Proposition~\ref{P-prop:decomp}, $4\eD(Q)=4\RD(\GQ)-2\dD(Q)$: at fixed Gibbs risk, each unit of disagreement removes two
units from the tandem bound.

\exo{5.6}
\statement{\textbf{(Application)} Three voters are evaluated on five equally likely examples $z_1,\dots,z_5$. Voter $h_1$ errs only on
$z_1$, $h_2$ errs on $z_2$ and $z_4$, and $h_3$ errs on $z_2$ and $z_3$; the posterior is $Q=(0.4,0.35,0.25)$. Compute $W_Q(z_j)$ and
$M_Q(z_j)$, the risk of the vote, the Gibbs risk, $\eD(Q)$ and $\dD(Q)$ (check the identity of Proposition~\ref{P-prop:decomp}), and the three
oracle bounds. Which one is the smallest, and is this consistent with Proposition~\ref{P-prop:compare}?}

\sol $W_Q(z_j)$ is the total weight of the voters that err on $z_j$:
\begin{center}\small
\begin{tabular}{lccccc}
\toprule
 & $z_1$ & $z_2$ & $z_3$ & $z_4$ & $z_5$ \\
\midrule
erring voters & $h_1$ & $h_2,h_3$ & $h_3$ & $h_2$ & none \\
$W_Q$ & $0.40$ & $0.60$ & $0.25$ & $0.35$ & $0$ \\
$M_Q=1-2W_Q$ & $0.2$ & $-0.2$ & $0.5$ & $0.3$ & $1$ \\
\bottomrule
\end{tabular}
\end{center}
The vote errs only on $z_2$, where $W_Q\geq\frac12$, so $\RD(\BQ)=0.2$. The Gibbs risk is $\E W_Q=\frac{1.6}5=0.32$ (a single voter errs on
$1.2$ examples out of $5$ on average). The joint error is $\eD(Q)=\E W_Q^2=\frac{0.16+0.36+0.0625+0.1225+0}{5}=0.141$, and the
disagreement is $\dD(Q)=\E[2W_Q(1-W_Q)]=\frac{2(0.24+0.24+0.1875+0.2275+0)}{5}=0.358$. It can be checked directly on pairs: $h_1$
disagrees with $h_2$ and with $h_3$ on $3$ examples out of $5$, and $h_2$ with $h_3$ on $2$, so
$\dD(Q)=2(0.4\cdot0.35\cdot0.6+0.4\cdot0.25\cdot0.6+0.35\cdot0.25\cdot0.4)=0.358$. The identity holds:
$\eD(Q)+\frac12\dD(Q)=0.141+0.179=0.32=\RD(\GQ)$. The oracle bounds are
\[
2\RD(\GQ)=0.64,\qquad 4\eD(Q)=0.564,\qquad
\mathcal C_Q=1-\frac{(1-0.64)^2}{1-0.716}=1-\frac{0.1296}{0.284}=0.544,
\]
where $1-2\dD(Q)=0.284$ is also $\E M_Q^2=\frac{0.04+0.04+0.25+0.09+1}5$. Since $\dD(Q)=0.358\geq\RD(\GQ)=0.32$, Proposition~\ref{P-prop:compare}
predicts $\mathcal C_Q\leq4\eD(Q)\leq2\RD(\GQ)$, which is what we find. All three are far from the true risk $0.2$: on five points, $W_Q$
has a large variance compared with its distance to $\frac12$.

\exo{5.7}
\statement{\textbf{(Application)} In Condorcet's model with $n=5$ independent voters of error $p=0.3$, compute the exact risk of the
vote, $\eD(Q)$, $\dD(Q)$ and the three oracle bounds. Check the regime predicted by Proposition~\ref{P-prop:compare}.}

\sol The vote errs when at least $3$ voters out of $5$ err:
\[
\RD(\BQ)=\binom53p^3(1-p)^2+\binom54p^4(1-p)+p^5=0.1323+0.02835+0.00243=0.1631.
\]
By Exercise~5.4, $\eD(Q)=p^2+\frac{p(1-p)}n=0.09+0.042=0.132$ and $\dD(Q)=2p(1-p)(1-\frac1n)=2\times0.21\times0.8=0.336$. The bounds are
$2p=0.6$, $4\eD(Q)=0.528$ and $\mathcal C_Q=1-\frac{(1-0.6)^2}{1-0.672}=1-\frac{0.16}{0.328}=0.512$. Since $\dD(Q)=0.336>p=0.3$, we are in the region where
the C-bound is the smallest, as observed. With $n=3$, $\dD(Q)=0.28<0.3$ and the first-order bound would be the best: $n=5$
is the smallest odd number of voters for which the C-bound wins, as stated in \cref{P-chap:mv}.

\exo{5.8}
\statement{\textbf{(Application)} A posterior with $\KL(Q\|P)=1$ has an empirical Gibbs risk $0.15$ on $m=1000$ labelled examples and
an empirical disagreement $0.25$ on $m'=1000$ examples. With $\delta=0.05$, compute $\overline r$, $\underline d$ and the PAC-Bayesian C-bound
\eqref{P-eq:cbound_pb}, the first-order certificate \eqref{P-eq:fo_pb} and the tandem certificate \eqref{P-eq:tnd_pb} (the empirical joint
error follows from Proposition~\ref{P-prop:decomp}, assuming the disagreement was measured on the same $1000$ labelled examples). Compare with the
value of the C-bound computed at the empirical values. What changes if $m'=10\,000$ unlabeled examples are available?}

\sol \emph{C-bound.} With the union bound, the confidence term is $\ln\frac{4\sqrt{1000}}{0.05}=7.836$. Hence
$\overline r=\klup\big(0.15,\frac{1+7.836}{1000}\big)=\klup(0.15,0.00884)=0.2014$ and
$\underline d=\kllow\big(0.25,\frac{2+7.836}{1000}\big)=\kllow(0.25,0.00984)=0.1928$, so that
\[
\RD(\BQ)\leq1-\frac{(1-2\times0.2014)^2}{1-2\times0.1928}=1-\frac{0.3565}{0.6144}=0.420.
\]
\emph{First order.} $2\klup\big(0.15,\frac{1+7.143}{1000}\big)=2\times0.1992=0.398$.

\emph{Tandem.} By Proposition~\ref{P-prop:decomp}, $\eS(Q)=0.15-\frac{0.25}2=0.025$, and $4\klup\big(0.025,\frac{2+7.143}{1000}\big)=0.209$.

\emph{Comparison.} At the empirical values, the C-bound would be $1-\frac{(1-0.3)^2}{1-0.5}=1-\frac{0.49}{0.5}=0.02$: a very diverse ensemble,
for which the oracle C-bound is excellent. The certificate, however, is $0.420$, worse than both other certificates. The
denominator $1-2d$ is small ($0.5$), so replacing $d=0.25$ by $\underline d=0.193$ changes it by $23\%$, and the numerator suffers from
$\overline r$ in the same way. With $m'=10\,000$ unlabeled examples, $\underline d=\kllow\big(0.25,\frac{2+\ln(4\sqrt{10^4}/0.05)}{10^4}\big)=0.2301$ and the
certificate becomes $0.340$: better, but still above the tandem certificate. Unlabeled data help the C-bound, but the
estimation of $\overline r$ remains a bottleneck.

\exo{5.9}
\statement{\textbf{(Proof)} For voters with values in $\{-1,+1\}$ and a fixed example $(x,y)$, show that
$\P_{(h,h')\sim Q^2}(h(x)\neq h'(x))=2W_Q(x,y)(1-W_Q(x,y))$. Deduce that $\E_\D M_Q^2=1-2\dD(Q)$ and $\RD(\GQ)=\eD(Q)+\frac12\dD(Q)$.}

\sol Fix $(x,y)$ and write $W=W_Q(x,y)$. Since the predictions take two values, $h(x)\neq h'(x)$ iff exactly one of the two voters
is wrong. As $h$ and $h'$ are drawn independently from $Q$, each errs with probability $W$, so
\begin{align*}
\P_{Q^2}(h(x)\neq h'(x))&=\P(h\text{ errs},h'\text{ correct})+\P(h\text{ correct},h'\text{ errs})\\
&=W(1-W)+(1-W)W=2W(1-W).
\end{align*}
Taking the expectation over $(x,y)\sim\D$ and exchanging the expectations (Fubini, non-negative integrand) gives
$\dD(Q)=\E_\D[2W_Q(1-W_Q)]$; similarly both voters err with probability $W^2$, whence $\eD(Q)=\E_\D W_Q^2$. By Exercise~1.9, $M_Q=1-2W_Q$,
so pointwise $M_Q^2=1-4W_Q+4W_Q^2=1-2\cdot2W_Q(1-W_Q)$, and $\E_\D M_Q^2=1-2\dD(Q)$. Finally, pointwise
$W_Q=W_Q^2+W_Q(1-W_Q)=W_Q^2+\frac12\cdot2W_Q(1-W_Q)$; taking expectations, $\RD(\GQ)=\eD(Q)+\frac12\dD(Q)$.

\emph{Remark.} The second moment of the margin is controlled by the disagreement alone: this is why the C-bound, which
uses $\E M_Q$ and $\E M_Q^2$, can be written with the Gibbs risk and the disagreement.

\exo{5.10}
\statement{\textbf{(Proof)} In Condorcet's model with $n=2k+1$ independent voters of error $p<\frac12$, use Hoeffding's inequality
(Lemma~\ref{P-lem:hoeffding}) to show that $\RD(\BQ)\leq e^{-2n(1/2-p)^2}$. Evaluate this bound for $p=0.3$ and $n\in\{5,25,101\}$, and compare it
with the exact risks and with the C-bound.}

\sol On a random example, let $Z_i=\ind{h_i(x)\neq y}$: the $Z_i$ are i.i.d.\ Bernoulli$(p)$ and $W_Q=\frac1n\sum_iZ_i$. The variables
$1-Z_i\in[0,1]$ are i.i.d.\ with mean $1-p$ and average $1-W_Q$, so Hoeffding's inequality applied to them gives, for $\varepsilon>0$,
$\P\big((1-p)-(1-W_Q)\geq\varepsilon\big)=\P(W_Q-p\geq\varepsilon)\leq e^{-2n\varepsilon^2}$. The vote errs iff $W_Q\geq\frac12$ (no tie is possible since $n$ is odd),
so with $\varepsilon=\frac12-p>0$,
\[
\RD(\BQ)=\P\big(W_Q-p\geq\tfrac12-p\big)\leq e^{-2n(1/2-p)^2}.
\]
For $p=0.3$, the exponent is $0.08\,n$:
\begin{center}\small
\begin{tabular}{cccc}
\toprule
$n$ & exact $\RD(\BQ)$ & Hoeffding $e^{-0.08n}$ & C-bound \\
\midrule
$5$ & $0.163$ & $0.670$ & $0.512$ \\
$25$ & $0.0175$ & $0.135$ & $0.174$ \\
$101$ & $1.29\times10^{-5}$ & $3.1\times10^{-4}$ & $0.049$ \\
\bottomrule
\end{tabular}
\end{center}
For large $n$, the Hoeffding bound decreases exponentially, like the true risk, whereas the C-bound only decreases like
$1/n$ (Exercise~5.4): the C-bound uses two moments of $W_Q$, Hoeffding the full independence of the voters. But independence
is exactly what fails for real ensembles and cannot be checked from data, whereas the C-bound only needs quantities that
PAC-Bayes can estimate. For small $n$ the C-bound is actually the better of the two.

%% ======================================================================
\section*{Section 6 -- Self-bounding learning}
\addcontentsline{toc}{section}{Section 6 -- Self-bounding learning}

\exo{6.1}
\statement{Reproduce the experiment of \cref{P-fig:uniformity} (draw $1000$ random classifiers and $m=100$ labels). How often
is the ``fixed-$h$'' Hoeffding bound violated? And the Occam bound? What happens to the Occam bound with $10^6$
classifiers?}

\sol See Exercise~3.4 for the code and the explanation. With $1000$ classifiers and $m=100$, the best empirical risk is
$0.34$ on average; the fixed-$h$ Hoeffding bound ($+0.122$) is violated in about $99.8\%$ of the runs, while the Occam
bound ($+\sqrt{(\ln1000+\ln20)/200}=+0.222$) is violated in about $0.3\%$ of them, below the nominal $5\%$. With $10^6$
classifiers, the best empirical risk drops to $0.26$ on average (a stronger selection), but the Occam price grows only
to $\sqrt{(\ln10^6+\ln20)/200}=0.290$: the certificate, about $0.55$ on average, is never violated in $200$ runs. It is
vacuous, which is the honest answer: with $100$ examples one cannot distinguish a million random classifiers from each
other. The logarithmic price keeps the bound valid for any number of candidates, at the cost of informativeness.

\exo{6.2}
\statement{Show that the value of $\lambda$ minimizing the Hoeffding-type bound \eqref{P-eq:catoni_hoeffding} depends on
$\KL(Q\|P)$, and explain why using it directly is forbidden, whereas using $\lambda^\star$ in the PAC-Bayes-$\lambda$ bound is allowed.}

\sol The bound is $\RS(\GQ)+\frac{A}{\lambda}+\frac\lambda8$ with $A=\frac{\KL(Q\|P)+\ln(1/\delta)}m$; its derivative in $\lambda$ vanishes
at $\lambda^\star=\sqrt{8A}=\sqrt{8(\KL(Q\|P)+\ln\frac1\delta)/m}$, which depends on $\KL(Q\|P)$, hence on the posterior and on the data.
In Corollary~\ref{P-cor:hoeffding_pb}, $\lambda$ is a parameter of the comparison function $\Delta(q,p)=\lambda(p-q)$: each value of
$\lambda$ defines a different event of probability at least $1-\delta$, and the corollary guarantees only that the event
associated with a $\lambda$ \emph{fixed in advance} holds. Choosing $\lambda$ after seeing the data means picking one event
among infinitely many, and the probability that the picked one fails is not controlled (the situation of
\cref{P-fig:uniformity}). It must be done on a finite grid with a union bound (Proposition~\ref{P-prop:union}).

In Corollary~\ref{P-cor:lambda}, on the contrary, the probabilistic event is the single event of Seeger's bound, which does
not involve $\lambda$; the PAC-Bayes-$\lambda$ inequality is derived from it by deterministic algebra valid for every
$\lambda\in(0,2)$ simultaneously. Plugging the data-dependent $\lambda^\star$ is therefore legitimate: the rule is ``free if it
indexes a family of inequalities all implied by the same event, paid if it changes the event''
(\cref{P-sec:sb_rules}).

\exo{6.3}
\statement{Using Proposition~\ref{P-prop:klgrad}, compute $\partial p/\partial q$ at $\hat q=0.1$, $\varepsilon=0.05$ (with $p\approx0.220$). Which
ingredient of the certificate is it more profitable to decrease: the empirical risk by $0.01$, or the KL by $5$ nats with
$m=1000$?}

\sol With $p=0.2201$, $q=0.1$: $\frac{p(1-p)}{p-q}=\frac{0.2201\times0.7799}{0.1201}=1.429$ and
$\ln\frac{p(1-q)}{q(1-p)}=\ln\frac{0.1981}{0.0780}=\ln2.540=0.932$. Hence
\[
\frac{\partial p}{\partial q}=1.429\times0.932=1.33,\qquad\frac{\partial p}{\partial\varepsilon}=1.43.
\]
Decreasing the empirical risk by $0.01$ decreases the certificate by about $1.33\times0.01=0.0133$ (the exact value is
$0.0135$). Decreasing the KL by $5$ nats decreases $\varepsilon$ by $5/1000=0.005$ and the certificate by about
$1.43\times0.005=0.0071$ (exact: $0.0073$). At this operating point, one point of empirical risk is worth about $9.3$ nats of
KL: it is more profitable to improve the fit. The balance changes with $m$: for $m=100$, the same $5$ nats would be
worth $0.071$, five times more than the gain in empirical risk. This is exactly the trade-off that the Gibbs
posterior or a gradient descent on the certificate resolves automatically.

\exo{6.4}
\statement{$(\star)$ Derive the ``quadratic'' relaxation of the table of \cref{P-sec:sb_optim} from the refined Pinsker
inequality. Writing $a=\sqrt{\varepsilon/2}$, show that it is smaller than McAllester's bound iff $2a+2\sqrt{\hat q+a^2}\leq1$; in
particular, for small $\varepsilon$ it is the case as soon as $\hat q<\frac14$.}

\sol Let $p=\klup(\hat q,\varepsilon)>\hat q$. By the refined Pinsker inequality, $\varepsilon\geq\kl(\hat q\|p)\geq\frac{(p-\hat q)^2}{2p}$, so
$(p-\hat q)^2\leq2\varepsilon p$, i.e.\ $p^2-2(\hat q+\varepsilon)p+\hat q^2\leq0$. The roots of this quadratic are
$\hat q+\varepsilon\pm\sqrt{2\hat q\varepsilon+\varepsilon^2}$, so
\[
p\leq\hat q+\varepsilon+\sqrt{2\hat q\varepsilon+\varepsilon^2}=\Big(\sqrt{\hat q+\tfrac\varepsilon2}+\sqrt{\tfrac\varepsilon2}\Big)^2,
\]
where the last equality is checked by expanding: $(\sqrt{\hat q+a^2}+a)^2=\hat q+2a^2+2a\sqrt{\hat q+a^2}$ with $2a^2=\varepsilon$ and
$2a\sqrt{\hat q+a^2}=\sqrt{2\varepsilon}\sqrt{\hat q+\varepsilon/2}=\sqrt{2\hat q\varepsilon+\varepsilon^2}$. McAllester's bound is $\hat q+a$. Then
\[
(\sqrt{\hat q+a^2}+a)^2\leq\hat q+a\iff2a^2+2a\sqrt{\hat q+a^2}\leq a\iff2a+2\sqrt{\hat q+a^2}\leq1
\]
(dividing by $a>0$). When $\varepsilon\to0$, $a\to0$ and the condition becomes $2\sqrt{\hat q}<1$, i.e.\ $\hat q<\frac14$. So the quadratic
bound improves on McAllester's for all classifiers better than $25\%$ error, provided the complexity term is small; for
$\hat q$ close to $\frac12$, Pinsker's inequality is tight and McAllester's bound is better than this relaxation (but never better
than the kl itself).

\exo{6.5}
\statement{\textbf{(Comprehension)} For each of the following practices, say whether it is free, requires a union bound, or is
forbidden, with a one-sentence justification. (a)~Stopping the gradient descent on the certificate when the certificate stops
decreasing. (b)~Choosing the scale $\sigma_0$ of a Gaussian prior among $5$ values after computing the five certificates.
(c)~Computing the certificate with $\delta=0.01$, reporting it with a $99\%$ claim if it is below a target value, and otherwise reporting the certificate computed with $\delta=0.05$ with a $95\%$ claim.
(d)~Choosing the temperature of a Gibbs posterior by cross-validation on the sample used in the bound.}

\sol (a)~\emph{Free.} The number of iterations only determines which posterior is output; the stopping rule is part of the
learning rule $S\mapsto Q_S$, and Theorem~\ref{P-thm:sb_valid} covers any stopping time.
(b)~\emph{Union bound.} The prior defines the probabilistic event of the bound; choosing among $5$ priors requires $\delta/5$ for
each, i.e.\ $\ln5\approx1.6$ extra nats.
(c)~\emph{Forbidden} as stated: the confidence level is the probability of an event fixed in advance, and a confidence level chosen after seeing the numbers has no frequentist meaning (the $99\%$ claim is made precisely on the samples where the certificate happens to be small). Note that the two events are nested, since the certificate computed with $\delta=0.01$ is always at least the one computed with $\delta=0.05$: reporting both certificates is valid, with the overall confidence $95\%$.
(d)~\emph{Free.} The temperature only indexes posteriors, all certified by the same event; any way of choosing it, including
cross-validation on the same sample, leaves the certificate valid (it may simply not be the smallest one).

\exo{6.6}
\statement{\textbf{(Comprehension)} True or false? Justify briefly. (a)~A self-bounding certificate is an unbiased estimate of the
true risk of the learned predictor. (b)~If the optimizer of a self-bounding algorithm is stuck in a poor local minimum, the
reported certificate may be wrong. (c)~A test-set bound remains valid if the test set is also used to choose among several
predictors.}

\sol (a)~\emph{False.} A certificate is a high-probability \emph{upper} bound, pessimistic by construction: in the running example,
the tandem certificate of the uniform vote is $0.387$ for a test risk of $0.044$. It should never be read as an estimate.
(b)~\emph{False.} The bound holds for all posteriors simultaneously, so the posterior found by a poor optimizer is certified as
well; its certificate is simply larger than it could have been.
(c)~\emph{False.} The test-set bound (Proposition~\ref{P-prop:testset}) requires the predictor to be independent of the test sample.
Selecting the best of $K$ predictors on the test set is a selection with the data, which must be paid for with $\delta/K$ (the
Occam bound with $K$ hypotheses); otherwise we are back to the situation of \cref{P-fig:uniformity}.

\exo{6.7}
\statement{\textbf{(Application)} In the model-selection example, the tandem certificates are computed with $\delta/6$. For a forest with an
OOB empirical joint error $\eS(Q)=0.05$, $\KL(Q\|P)=2$ and $n_0=423$, compute the tandem certificate \eqref{P-eq:tnd_pb} with $\delta=0.05$,
with $\delta/6$, and with $\delta/100$. With a countable family and $\pi_k=\frac6{\pi^2k^2}$, what is the extra cost, in nats, for the
candidates $k=1$ and $k=10$?}

\sol The budget is $\varepsilon=\frac{2\KL+\ln(2\sqrt{n_0}/\delta)}{n_0}$ with $\ln\frac{2\sqrt{423}}{0.05}=6.713$.
\begin{itemize}
\item $\delta=0.05$: $\varepsilon=\frac{4+6.713}{423}=0.02533$ and $4\klup(0.05,0.02533)=0.458$;
\item $\delta/6$: $\varepsilon$ increases by $\frac{\ln6}{423}=0.00424$, to $0.02956$, and the certificate becomes $0.484$;
\item $\delta/100$: $\varepsilon=0.03621$ and the certificate becomes $0.522$.
\end{itemize}
Choosing among $6$ candidates costs $1.79$ nats and about $0.026$ of certificate, among $100$ candidates about $0.065$.
For the countable family, candidate $k$ is certified with $\pi_k\delta$, i.e.\ at the extra cost $\ln\frac1{\pi_k}=2\ln k+\ln\frac{\pi^2}6$ nats:
$0.50$ nats for $k=1$ and $2\ln10+0.50=5.10$ nats for $k=10$. The ordering of the candidates must be fixed in advance:
candidates listed first are cheaper.

\exo{6.8}
\statement{\textbf{(Application)} With $m=1000$ examples, the hold-out approach trains a forest on $700$ examples and its vote errs on $30$ of
the $300$ remaining ones. The self-bounding approach trains a forest on the $1000$ examples and obtains an OOB empirical joint
error of $0.06$ with $\KL(Q\|P)=1.5$ and $n_0=370$. With $\delta=0.05$, compute the test-set bound and the tandem certificate. Below
which value of the empirical joint error would the tandem certificate beat the hold-out one? What is the tandem certificate
when $\eS(Q)=0$?}

\sol \emph{Hold-out.} $\widehat R_T=\frac{30}{300}=0.10$ and $\klup\big(0.10,\frac{\ln20}{300}\big)=\klup(0.10,0.00999)=0.148$.

\emph{Self-bounding.} $\varepsilon=\frac{2\times1.5+\ln(2\sqrt{370}/0.05)}{370}=\frac{3+6.646}{370}=0.02607$ and $4\klup(0.06,0.02607)=0.518$.

\emph{Threshold.} The tandem certificate $4\klup(\hat e,0.02607)$ is increasing in $\hat e$; solving
$4\klup(\hat e,0.02607)=0.148$ by bisection gives $\hat e\approx0.0034$. Even with no joint error at all,
$4\klup(0,\varepsilon)=4(1-e^{-\varepsilon})=0.103$: the factor $4$ applied to the complexity term alone already costs $0.10$.

\emph{Comment.} This reproduces the conclusion of \cref{P-fig:sb_holdout}: for majority votes, the factor $4$ and the doubled KL of
the tandem bound cost more than the $300$ examples withheld by the hold-out. The self-bounding approach wins on the
quality of the vote (all data are used to learn) and on model selection, not on the tightness of the certificate.

\exo{6.9}
\statement{\textbf{(Proof)} Prove Proposition~\ref{P-prop:union} in detail: write the failure event of each bound, apply the union bound, and
explain why the certificate $\min_k\mathcal B_k(Q_S,S,\pi_k\delta)$ is valid for any learning rule and any data-dependent choice of $k$.
Check that $\pi_k=\frac{6}{\pi^2k^2}$, $k\geq1$, is a valid choice of weights.}

\sol For each $k$, let $F_k=\{S:\exists Q,\ \RD(Q)>\mathcal B_k(Q,S,\pi_k\delta)\}$ be the failure event of the $k$-th bound, used with
confidence $\pi_k\delta$. By assumption $\P(F_k)\leq\pi_k\delta$. By the union bound (subadditivity of probabilities, valid for a finite
or countable family),
\[
\P\Big(\bigcup_kF_k\Big)\leq\sum_k\P(F_k)\leq\sum_k\pi_k\delta=\delta.
\]
On the complement $G=\bigcap_kF_k^c$, which has probability at least $1-\delta$, we have $\RD(Q)\leq\mathcal B_k(Q,S,\pi_k\delta)$ for \emph{all} $k$ and
\emph{all} $Q$ simultaneously. For $S\in G$, take any posterior $Q_S$ and any index $k_S$, both possibly depending on $S$: the
inequality holds in particular for $(k_S,Q_S)$, and taking the minimum over $k$ of valid upper bounds gives an upper bound,
$\RD(Q_S)\leq\min_k\mathcal B_k(Q_S,S,\pi_k\delta)$. Since $G$ does not depend on the choices, the argument of Theorem~\ref{P-thm:sb_valid}
applies. Finally, $\pi_k>0$ and $\sum_{k\geq1}\frac6{\pi^2k^2}=\frac6{\pi^2}\cdot\frac{\pi^2}6=1$ (Basel problem), so these weights are admissible.

\exo{6.10}
\statement{\textbf{(Proof)} Let $\rho=\mathrm{softmax}(\theta)$ with $\theta\in\R^n$. Show that $\frac{\partial\rho_s}{\partial\theta_t}=\rho_s(\ind{s=t}-\rho_t)$, deduce the formula
\eqref{P-eq:softmax_grad} for $\nabla_\theta\mathcal B$, and show that the components of $\nabla_\theta\mathcal B$ sum to zero. Interpret this last fact.}

\sol Let $Z=\sum_ue^{\theta_u}$, so that $\rho_s=e^{\theta_s}/Z$ and $\partial Z/\partial\theta_t=e^{\theta_t}$. By the quotient rule,
\[
\frac{\partial\rho_s}{\partial\theta_t}=\frac{\ind{s=t}e^{\theta_s}Z-e^{\theta_s}e^{\theta_t}}{Z^2}=\rho_s\ind{s=t}-\rho_s\rho_t=\rho_s\big(\ind{s=t}-\rho_t\big).
\]
With $g=\nabla_\rho\mathcal B$, the chain rule gives
$\frac{\partial\mathcal B}{\partial\theta_t}=\sum_sg_s\rho_s(\ind{s=t}-\rho_t)=\rho_tg_t-\rho_t\sum_sg_s\rho_s=\rho_t\big(g_t-\rho^\top g\big)$, i.e.\
$\nabla_\theta\mathcal B=\rho\odot(g-(\rho^\top g)\mathbf1)$, which is \eqref{P-eq:softmax_grad}. Summing the components,
$\sum_t\rho_tg_t-(\rho^\top g)\sum_t\rho_t=\rho^\top g-\rho^\top g=0$.

\emph{Interpretation.} Adding the same constant $c$ to all logits does not change $\rho$ ($e^{\theta_s+c}/\sum_ue^{\theta_u+c}=\rho_s$), so
$\mathcal B$ is constant along the direction $\mathbf1$ and its gradient is orthogonal to $\mathbf1$. In particular, the constant $\mathbf1$ in
$\nabla_\rho\KL(\rho\|\pi)=\ln\frac\rho\pi+\mathbf1$ disappears after the projection: only relative weights matter.

%% ======================================================================
\section*{Section 7 -- Self-bounding algorithms for ensembles and beyond}
\addcontentsline{toc}{section}{Section 7 -- Self-bounding algorithms for ensembles and beyond}

\exo{7.1}
\statement{Derive the gradient \eqref{P-eq:grad_tnd_lambda} of the tandem objective, and check it numerically with finite
differences (notebook~4).}

\sol Write $c_1=\frac{4}{1-\lambda/2}$ and $c_2=\frac{4}{\lambda(1-\lambda/2)n_0}$, so that
$\mathcal B_{\mathrm{TND}}(\rho)=c_1\,\rho^\top\widehat{\mathbf L}\rho+c_2\big(2\KL(\rho\|\pi)+\ln\frac{2\sqrt{n_0}}\delta\big)$. The matrix $\widehat{\mathbf L}$ of
tandem losses is symmetric ($\widehat L_{tt'}=\widehat L_{t't}$), so $\nabla_\rho(\rho^\top\widehat{\mathbf L}\rho)=2\widehat{\mathbf L}\rho$. For the KL,
$\KL(\rho\|\pi)=\sum_t\rho_t\ln\rho_t-\rho_t\ln\pi_t$, and $\partial_{\rho_t}=\ln\rho_t+1-\ln\pi_t$, i.e.\
$\nabla_\rho\KL=\ln\frac\rho\pi+\mathbf1$ (componentwise). Hence
\[
\nabla_\rho\mathcal B_{\mathrm{TND}}=4\Big[\frac{2\widehat{\mathbf L}\rho}{1-\lambda/2}+\frac{2(\ln\frac\rho\pi+\mathbf 1)}{\lambda(1-\lambda/2)n_0}\Big].
\]
With the softmax parametrization $\rho=\mathrm{softmax}(\theta)$, the Jacobian is $\frac{\partial\rho_s}{\partial\theta_t}=\rho_s(\ind{s=t}-\rho_t)$,
so $\nabla_\theta\mathcal B=\rho\odot(g-(\rho^\top g)\mathbf 1)$ with $g=\nabla_\rho\mathcal B$, which is \eqref{P-eq:softmax_grad}. (The constant $\mathbf1$
in the KL gradient disappears in this projection, as it should: it corresponds to the direction that changes
$\sum_t\rho_t$.) \emph{Numerical check.} On the forest of the running example ($n_0=423$), at random logits $\theta$ and
$\lambda=0.7$, the centred finite differences with step $10^{-6}$ agree with the analytical gradient up to a relative error
of $6\times10^{-8}$:
\begin{verbatim}
num = np.array([(f(th + 1e-6*e) - f(th - 1e-6*e)) / 2e-6
                for e in np.eye(100)])
print(np.linalg.norm(g - num) / np.linalg.norm(num))   # 5.6e-08
\end{verbatim}

\exo{7.2}
\statement{Explain why the uniform prior over the trees of a random forest would \emph{not} be a valid prior if the empirical
losses were computed on the whole training sample.}

\sol The uniform prior on the \emph{indices} of the trees is fixed in advance, and in this sense it is legitimate.
The problem lies in the third step of the proof of Theorem~\ref{P-thm:general}: for each index $t$, we need the empirical loss
$\widehat L_t$ to be an average of i.i.d.\ variables of mean $\RD(h_t)$, independent of $h_t$. If $\widehat L_t$ is computed on the whole
training sample, it includes the examples used to grow $h_t$, on which a fully grown tree makes almost no error: $\widehat L_t$ is
then a strongly biased estimate of $\RD(h_t)$ (close to $0$ for unpruned trees), and Maurer's lemma does not apply.
Equivalently, the prior on the trees $h_t$ themselves (not on their indices) depends on the sample used to evaluate
them. The certificate would be meaningless: it would certify a near-zero risk for a forest whose trees err on $19\%$ of
the test digits. The out-of-bag construction restores the independence by evaluating each tree only on the examples
it has not seen (Theorem~\ref{P-thm:oob}).

\exo{7.3}
\statement{In PBGD, show that when all the training points have a large normalized margin, the gradient of the loss term
vanishes, and interpret the resulting value of $\|\mathbf w\|$.}

\sol The gradient of the objective is $\nabla F(\mathbf w)=-C\sum_i\varphi(t_i)\frac{y_ix_i}{\|x_i\|}+\mathbf w$ with
$t_i=\frac{y_i\langle\mathbf w,x_i\rangle}{\|x_i\|}$ and $\varphi$ the standard Gaussian density. Since $\|\frac{y_ix_i}{\|x_i\|}\|=1$, the loss
part has norm at most $C\sum_i\varphi(t_i)\leq Cm\,\varphi(t_{\min})$, and $\varphi(t)=\frac{e^{-t^2/2}}{\sqrt{2\pi}}\to0$ very fast: if all
$t_i\geq t_{\min}$ large, the loss gradient vanishes. Then $\nabla F(\mathbf w)\approx\mathbf w\neq0$, and gradient descent shrinks
$\mathbf w$. Shrinking $\mathbf w$ decreases all normalized margins proportionally ($t_i$ is linear in $\|\mathbf w\|$), until the smallest
ones become of order $1$--$2$, where $\varphi(t_i)$ is no longer negligible. At a stationary point,
$\mathbf w=C\sum_i\varphi(t_i)\frac{y_ix_i}{\|x_i\|}$: $\mathbf w$ is a combination of the points with moderate margins, the analogue
of support vectors. If the data are separable with geometric normalized margin $\gamma$ in direction $\mathbf u$, the
stationary $\mathbf w=s\mathbf u$ has $s\gamma=O(1)$, i.e.\ $\|\mathbf w\|\approx\frac{c}{\gamma}$ with $c$ growing slowly (logarithmically) with $C$, since $\varphi(t)$ decays like
$e^{-t^2/2}$. The KL, $\frac12\|\mathbf w\|^2\approx\frac{c^2}{2\gamma^2}$, is thus inversely proportional to the squared margin: large margins
give small certificates, which is the PAC-Bayesian explanation of the SVM. On the example of \cref{P-fig:pbgd}, the
norm grows from $3.5$ ($C=1$) to $13.6$ ($C=100$) only because of the noisy points, which pull $\mathbf w$ as long as their
margin is moderate.

\exo{7.4}
\statement{With notebook~4, compute the test disagreement of the C-bound posterior of \cref{P-tab:rf_posteriors} and check
whether the condition $d\geq R$ of Proposition~\ref{P-prop:compare} holds. Why is its certificate nevertheless larger than the
tandem certificate?}

\sol For the posterior minimizing the PAC-Bayesian C-bound, the test Gibbs risk is $R=0.176$ and the test
disagreement is $d=0.246$: the condition $d\geq R$ holds, and indeed the \emph{oracle} C-bound is $0.171$, smaller than the
oracle tandem bound $4e=0.210$. Its certificate, however, is $0.531$, against $0.361$ for the TND-optimized posterior. The
reason is statistical. The certificate replaces $R$ by an upper confidence bound and $d$ by a lower confidence bound,
each with $\delta/2$: from the OOB estimates $\hat R=0.153$ ($n_1=747$) and $\hat d=0.224$ ($n_3=423$ pairs), with $\KL=0.73$,
we get $\overline r=0.212$ and $\underline d=0.147$. The disagreement is estimated on the smaller pairwise sample with a doubled KL,
and it enters through $\frac1{1-2d}$, which amplifies its estimation error: plugging $\overline r$ and $\underline d$ gives
$1-\frac{(1-0.424)^2}{1-0.294}=0.531$. The tandem certificate needs a single kl inversion of a small quantity ($\hat e=0.041$),
where the kl is at its best. This is the general lesson of Proposition~\ref{P-prop:compare}: the best oracle bound is not
necessarily the best certificate; the number of quantities to estimate, and their size, matter.

\exo{7.5}
\statement{$(\star)$ For a stochastic network with $D=10^5$ weights, $\sigma_j=\sigma_0=0.03$ and $\|\boldsymbol\mu-\boldsymbol\mu_0\|=3$, compute
$\KL(Q\|P)$. How many examples are needed for the complexity term $\KL/m$ to be below $0.01$? What does this suggest
about data-dependent priors?}

\sol With $\sigma_j=\sigma_0$, the terms $\frac{\sigma_j^2}{\sigma_0^2}-1+\ln\frac{\sigma_0^2}{\sigma_j^2}$ vanish, and
\[
\KL(Q\|P)=\frac12\sum_j\frac{(\mu_j-\mu_{0,j})^2}{\sigma_0^2}=\frac{\|\boldsymbol\mu-\boldsymbol\mu_0\|^2}{2\sigma_0^2}=\frac{9}{2\times0.0009}=5000\text{ nats}.
\]
The complexity term $\KL/m$ is below $0.01$ only when $m\geq5000/0.01=500\,000$ examples, about eight times the size of
MNIST's training set, and this ignores the confidence term and the kl inversion. Note that the dimension $D$ does not
appear: what matters is how far, in units of $\sigma_0$, the posterior mean has travelled from the prior mean. Moving
$3$ units of Euclidean distance with a noise scale of $0.03$ is enormous. This suggests two levers: take $\sigma_0$ as large
as the network tolerates (chosen on a grid with a union bound), and above all bring $\boldsymbol\mu_0$ close to $\boldsymbol\mu$ with a
data-dependent prior: if $\boldsymbol\mu_0$ is a network trained on half of the data, the final network may for instance move by a
distance of $0.3$ instead of $3$, which divides the KL by $100$ ($50$ nats): the certificate then becomes informative
with a few tens of thousands of examples. This is the approach of \citet{dziugaite2018data} and \citet{perez2021tighter}
(Proposition~\ref{P-prop:ddprior}).

\exo{7.6}
\statement{\textbf{(Comprehension)} True or false? Justify briefly. (a)~The OOB bound is invalid because every training example has
been used to build some of the trees. (b)~Adding trees to a forest with the uniform posterior increases $\KL(\rho\|\pi)$. (c)~In
PBGD, the parameter $C$ must be selected by cross-validation, since the bound cannot be used to choose it. (d)~The
first-order posterior of \cref{P-tab:rf_posteriors} has a small Gibbs risk, hence necessarily a good vote.}

\sol (a)~\emph{False.} What matters is that each tree is evaluated only on examples it has not seen: conditionally on its
bootstrap sample, $h_t$ is fixed and its OOB sample is i.i.d.\ and independent of it (Theorem~\ref{P-thm:oob}). An example may help
build $h_1$ and certify $h_2$.
(b)~\emph{False.} With $\rho=\pi$ uniform, $\KL(\rho\|\pi)=0$ whatever the number of trees. What changes is $n_0$, a minimum over
more trees (or pairs of trees), which decreases slightly.
(c)~\emph{False.} $C$ only indexes the posteriors $Q_{\mathbf w}$, all compared with the same prior in the same bound: it is free,
and PBGD can select it by minimizing its own certificate (\cref{P-fig:pbgd}).
(d)~\emph{False.} Its Gibbs risk is the smallest ($0.158$), but it keeps two effective trees and its vote errs on $15\%$ of the
test digits, against $4.4\%$ for the uniform forest: a small Gibbs risk without diversity does not make a good vote.

\exo{7.7}
\statement{\textbf{(Application)} For the \texttt{digits} forest ($m=1257$), compute the expected fraction and number of examples that are
out-of-bag for one tree, and for two given trees, when the bootstrap samples have size $m$ and when they have size
$\lfloor m/2\rfloor=628$. Compare with the values $n_1=747$ and $n_2=423$ used in the notes, and explain the difference.}

\sol A bootstrap sample of size $k$ draws $k$ indices uniformly with replacement, so a given example is out-of-bag with
probability $(1-\frac1m)^k$, and out-of-bag for two independent bootstrap samples with probability $(1-\frac1m)^{2k}$. With $m=1257$:
\begin{center}\small
\begin{tabular}{ccccc}
\toprule
bootstrap size $k$ & OOB, one tree & expected size & OOB, two trees & expected size \\
\midrule
$1257$ & $0.3677\approx e^{-1}$ & $462.2$ & $0.1352\approx e^{-2}$ & $170.0$ \\
$628$ & $0.6067\approx e^{-1/2}$ & $762.6$ & $0.3680\approx e^{-1}$ & $462.6$ \\
\bottomrule
\end{tabular}
\end{center}
Halving the bootstrap size multiplies the sample available for the tandem bound by $2.7$. The values of the notes, $n_1=747$
and $n_2=423$, are slightly below the expected sizes $762.6$ and $462.6$ because Theorem~\ref{P-thm:oob} uses the \emph{smallest}
OOB set over the $100$ trees, and the smallest intersection over the $4950$ pairs. A simulation of $100$ bootstrap
samples of size $628$ (20 random seeds) gives a smallest OOB set of $743$ on average (between $733$ and $748$) and a smallest
pairwise intersection of $423$ on average (between $416$ and $428$), in agreement with the notes.

\exo{7.8}
\statement{\textbf{(Application)} Run \cref{P-algo:fo} by hand on $n=4$ trees with OOB losses $\widehat L=(0.15,0.17,0.20,0.25)$, uniform prior,
$n_0=700$ and $\delta=0.05$: compute $\rho$ for $\lambda=1$, the updated $\lambda$, and $\rho$ again. Compare the final first-order certificate with
the one of the uniform posterior, and comment on the weights.}

\sol \emph{First step, $\lambda=1$.} $\rho_t\propto e^{-700\widehat L_t}$; relative to the first tree, the weights are $1$, $e^{-14}=8.3\times10^{-7}$,
$e^{-35}=6.3\times10^{-16}$ and $e^{-70}$: $\rho\approx(1,\ 8.3\times10^{-7},\ 0,\ 0)$, with $\rho^\top\widehat L=0.1500$ and
$\KL(\rho\|\pi)\approx\ln4=1.386$. \emph{Update of $\lambda$.} With $\ln\frac{2\sqrt{700}}{0.05}=6.964$,
\[
\lambda=\frac{2}{\sqrt{\frac{2\times700\times0.15}{1.386+6.964}+1}+1}=\frac{2}{\sqrt{26.15}+1}=0.327.
\]
\emph{Second step.} $\rho_t\propto e^{-0.327\times700\,\widehat L_t}=e^{-229\,\widehat L_t}$; relative weights $1$, $e^{-4.58}=0.0102$,
$e^{-11.45}=1.1\times10^{-5}$ and $e^{-22.9}=1.1\times10^{-10}$: $\rho=(0.9898,\ 0.0102,\ 0.00001,\ 0.0000)$, with $\rho^\top\widehat L=0.1502$ and $\KL=1.329$. The next
update gives $\lambda=0.3261$ and the iterations are stable from then on (\texttt{optimize\_fo} of \texttt{pbtools}). The
PAC-Bayes-$\lambda$ bound on the Gibbs risk is $0.223$, and the first-order certificate on the vote is
$2\klup\big(0.1502,\frac{1.329+6.964}{700}\big)=2\times0.2105=0.421$. The uniform posterior has $\rho^\top\widehat L=0.1925$, $\KL=0$ and the certificate
$2\klup(0.1925,\frac{6.964}{700})=2\times0.2520=0.504$.

\emph{Comment.} The certificate improves by $0.08$, but the posterior is essentially a Dirac on the first tree: the ``vote'' is
a single tree. With $n_0=700$, a difference of $0.02$ in OOB loss is $14$ errors, which the Gibbs posterior turns into a
factor $e^{-14\lambda}$. This is the mechanism by which first-order minimization destroys diversity (\cref{P-fig:rf_weights}).

\exo{7.9}
\statement{\textbf{(Application)} For a Gaussian posterior $Q_{\mathbf w}=\mathcal N(\mathbf w,I_d)$ with $\|\mathbf w\|=2$, three training points have
normalized margins $y_i\langle\mathbf w,x_i\rangle/\|x_i\|$ equal to $2$, $0.5$ and $-1$. Compute the empirical Gibbs risk \eqref{P-eq:probit}, the
training error of the vote and $\KL(Q_{\mathbf w}\|P)$. Redo the computation for $3\mathbf w$, and relate the result to \cref{P-fig:linear}.}

\sol By \eqref{P-eq:probit}, each point contributes $\Phi(-t_i)$ where $t_i$ is its normalized margin:
$\Phi(-2)=0.0228$, $\Phi(-0.5)=0.3085$ and $\Phi(1)=0.8413$, so $\RS(G_{Q_{\mathbf w}})=\frac{0.0228+0.3085+0.8413}3=0.391$. The vote is $\sign\langle\mathbf w,x\rangle$,
which errs only on the point with a negative margin: training error $\frac13$. The KL is $\frac12\|\mathbf w\|^2=2$. For $3\mathbf w$, the margins
become $6$, $1.5$ and $-3$: $\Phi(-6)=10^{-9}$, $\Phi(-1.5)=0.0668$ and $\Phi(3)=0.9987$, so the Gibbs risk is $0.355$; the vote and its
training error ($\frac13$) are unchanged, and the KL becomes $\frac12\times36=18$.

\emph{Comment.} Scaling $\mathbf w$ does not change the vote but makes the sampled classifiers agree with it: the Gibbs risk
decreases towards the training error of the vote ($\frac13$, its limit as the scale goes to infinity) while the KL grows
quadratically. This is the trade-off of \cref{P-fig:linear}, which PBGD resolves by choosing the norm of $\mathbf w$.

\exo{7.10}
\statement{\textbf{(Proof)} Let $Q_{\mathbf w}=\mathcal N(\mathbf w,I_d)$ over linear voters $h_{\mathbf v}(x)=\sign\langle\mathbf v,x\rangle$, and $x\neq0$. Show that
$\P_{\mathbf v\sim Q_{\mathbf w}}(y\langle\mathbf v,x\rangle\leq0)=\Phi(-y\langle\mathbf w,x\rangle/\|x\|)$, and deduce that the majority vote is $B_{Q_{\mathbf w}}(x)=\sign\langle\mathbf w,x\rangle$
(up to the tie $\langle\mathbf w,x\rangle=0$) and the formula \eqref{P-eq:probit}.}

\sol Write $\mathbf v=\mathbf w+\boldsymbol\xi$ with $\boldsymbol\xi\sim\mathcal N(\mathbf 0,I_d)$. Then $y\langle\mathbf v,x\rangle=y\langle\mathbf w,x\rangle+y\langle\boldsymbol\xi,x\rangle$, and $y\langle\boldsymbol\xi,x\rangle$ is a
linear combination of independent standard Gaussian variables, hence Gaussian with mean $0$ and variance $y^2\|x\|^2=\|x\|^2$. So
$y\langle\boldsymbol\xi,x\rangle=\|x\|G$ with $G\sim\mathcal N(0,1)$ and
\[
\P\big(y\langle\mathbf v,x\rangle\leq0\big)=\P\Big(G\leq-\frac{y\langle\mathbf w,x\rangle}{\|x\|}\Big)=\Phi\Big(-\frac{y\langle\mathbf w,x\rangle}{\|x\|}\Big).
\]
Since $\P(\langle\mathbf v,x\rangle=0)=0$, the event $\{y\langle\mathbf v,x\rangle\leq0\}$ is, almost surely, the event that $h_{\mathbf v}$ errs, so this is
$W_{Q_{\mathbf w}}(x,y)$. The vote errs iff $W_{Q_{\mathbf w}}\geq\frac12$ (Exercise~1.9), i.e.\ iff $y\langle\mathbf w,x\rangle\leq0$, because $\Phi(-t)\geq\frac12\iff t\leq0$.
Hence the vote agrees with $\sign\langle\mathbf w,x\rangle$ except possibly on the tie $\langle\mathbf w,x\rangle=0$. Finally, averaging
$W_{Q_{\mathbf w}}(x_i,y_i)$ over the sample gives $\RS(G_{Q_{\mathbf w}})=\frac1m\sum_i\Phi\big(-y_i\langle\mathbf w,x_i\rangle/\|x_i\|\big)$, which is
\eqref{P-eq:probit}.

%% ======================================================================
\section*{Section 8 -- Further topics and recent advances}
\addcontentsline{toc}{section}{Section 8 -- Further topics and recent advances}

\exo{8.1}
\statement{\textbf{(Comprehension)} A practitioner collects examples one at a time, recomputes Seeger's bound with $\delta=0.05$ after
each new example, and stops as soon as the certificate is below a target. Explain why the reported certificate is not
guaranteed to hold with probability $0.95$. Show that applying the bound for sample size $m$ with confidence
$\delta_m=\frac{\delta}{m(m+1)}$ gives a bound valid simultaneously for all $m\geq8$, and compute the extra cost in nats and the resulting
Seeger bound for $m=1000$, $\RS(\GQ)=0.1$, $\KL(Q\|P)=5$, $\delta=0.05$.}

\sol For each \emph{fixed} $m$, the bound fails with probability at most $\delta$. But the stopping time $\tau$ (the first $m$ at
which the certificate is below the target) depends on the data, and the certificate is reported precisely when it
happens to be small. The event ``the bound fails at $\tau$'' is contained in the union of the failure events over all
$m$, whose probability can be much larger than $\delta$: this is the selection bias of \cref{P-fig:uniformity}, now over sample
sizes instead of classifiers.

The fix is a union bound over $m$. Since $\frac1{m(m+1)}=\frac1m-\frac1{m+1}$, the sum telescopes: $\sum_{m\geq1}\frac{\delta}{m(m+1)}=\delta$. Applying
Seeger's bound at each $m\geq8$ with $\delta_m$, the probability that at least one of them fails is at most
$\sum_{m\geq8}\delta_m\leq\delta$; on the complementary event, the bound holds for every $m$ and every $Q$ at once, hence at any
stopping time. The price is $\ln\frac1{\delta_m}-\ln\frac1\delta=\ln\big(m(m+1)\big)$ extra nats: for $m=1000$, $\ln(1\,001\,000)=13.8$ nats, to
be compared with $\ln\frac{2\sqrt{1000}}{0.05}=7.14$. Seeger's bound goes from $\klup\big(0.1,\frac{5+7.14}{1000}\big)=0.153$ to
$\klup\big(0.1,\frac{5+7.14+13.82}{1000}\big)=0.182$.

\emph{Remark.} The price grows only logarithmically with $m$. The supermartingale techniques mentioned in the section are a more
refined way to obtain such anytime-valid guarantees.

\exo{8.2}
\statement{\textbf{(Proof)} Let $\Hcal$ be finite, $S\mapsto Q_S$ a learning rule, and write $\bar Q=\E_S Q_S$ for the average posterior. Show
that, for any prior $P$, $\E_S\KL(Q_S\|P)=\E_S\KL(Q_S\|\bar Q)+\KL(\bar Q\|P)$. Deduce that $\bar Q$ minimizes $\E_S\KL(Q_S\|P)$ over
data-independent priors, and that the minimum equals the mutual information $I(S;h)$ between the sample and a hypothesis
$h\sim Q_S$.}

\sol If $\bar Q(h)=0$, then $Q_S(h)=0$ for almost every $S$, so all the terms below are well defined (if $\KL(\bar Q\|P)=+\infty$, both sides
are infinite). For every $S$, splitting the logarithm,
\[
\KL(Q_S\|P)=\sum_hQ_S(h)\ln\frac{Q_S(h)}{\bar Q(h)}+\sum_hQ_S(h)\ln\frac{\bar Q(h)}{P(h)}=\KL(Q_S\|\bar Q)+\sum_hQ_S(h)\ln\frac{\bar Q(h)}{P(h)}.
\]
The last sum is linear in $Q_S$; its expectation over $S$ is $\sum_h\bar Q(h)\ln\frac{\bar Q(h)}{P(h)}=\KL(\bar Q\|P)$. Taking expectations gives the
identity. Since $\KL(\bar Q\|P)\geq0$ with equality iff $P=\bar Q$, the prior $P=\bar Q$ minimizes $\E_S\KL(Q_S\|P)$, and the minimum is
$\E_S\KL(Q_S\|\bar Q)$. Now consider the joint distribution of $(S,h)$ where $S\sim\D^m$ and $h\,|\,S\sim Q_S$: the marginal law of $h$ is
$\bar Q$, and the mutual information is by definition the expected divergence between the conditional and the marginal
laws of $h$, $I(S;h)=\E_S\KL(Q_S\|\bar Q)$.

\emph{Remark.} $\bar Q$ depends on $\D$ but not on the sample: it is a legitimate, although unknown, prior. Averaging a
PAC-Bayesian bound over $S$ with this prior replaces the KL term by $I(S;h)$, which is the bridge with the
information-theoretic bounds of \citet{xu2017information}.

\exo{8.3}
\statement{\textbf{(Comprehension)} Explain why the disagreement $\dD(Q)$ of a vote can be estimated on unlabeled target examples, whereas
its Gibbs risk and its joint error cannot. Suppose that the disagreement measured on unlabeled target data is $0.30$ and that
the target joint error is (somehow) known to be $0.05$: compute the target Gibbs risk and the oracle C-bound on the target
risk of the vote.}

\sol The disagreement $\dD(Q)=\E_{(h,h')\sim Q^2}\P_{x}(h(x)\neq h'(x))$ only involves the predictions of the voters on the inputs $x$; the
labels do not appear. It can therefore be estimated on target inputs alone. The Gibbs risk (``$h$ errs'') and the joint error
(``$h$ and $h'$ both err'') compare predictions with the label $y$, which is not observed on the target domain. By
Proposition~\ref{P-prop:decomp}, the target Gibbs risk is $\eD(Q)+\frac12\dD(Q)=0.05+0.15=0.20$, and the oracle C-bound is
$1-\frac{(1-0.4)^2}{1-0.6}=1-\frac{0.36}{0.40}=0.10$. In practice the joint error on the target is unknown; the domain adaptation bounds of
\citet{germain2013pac,germain2020pac} relate the unknown target quantities to source quantities and to a divergence between the
domains built on the disagreement, which is precisely the quantity that unlabeled target data give access to.

\exo{8.4}
\statement{\textbf{(Comprehension)} A certificate $\mathcal B$ holds for the Gibbs risk $\RD(\GQ)$ with probability at least $1-\delta$. A single
network $h\sim Q$ is drawn once and deployed. Is $\RD(h)\leq\mathcal B$ guaranteed with high probability? Using Markov's inequality over
$h\sim Q$, show that $\RD(h)\leq\mathcal B/\eta$ with probability at least $1-\delta-\eta$, and compute this guarantee for $\mathcal B=0.05$ and
$\eta=0.1$. What do disintegrated bounds improve?}

\sol No: $\mathcal B$ bounds the \emph{average} $\E_{h\sim Q}\RD(h)$, and a single draw may be worse than average (for instance, if $Q$
puts mass $0.9$ on voters of risk $0$ and $0.1$ on voters of risk $0.5$, the average is $0.05$, yet one draw in ten has risk $0.5$).
On the event of probability at least $1-\delta$ over $S$ where $\E_Q\RD(h)\leq\mathcal B$, Markov's inequality over $h\sim Q$ gives
$\P_{h\sim Q}(\RD(h)\geq\mathcal B/\eta)\leq\frac{\eta\,\E_Q\RD(h)}{\mathcal B}\leq\eta$. The joint draw of $(S,h)$ fails if either the certificate fails (probability
at most $\delta$) or the draw is bad (probability at most $\eta$), so $\RD(h)\leq\mathcal B/\eta$ with probability at least $1-\delta-\eta$. With
$\mathcal B=0.05$ and $\eta=0.1$: $\RD(h)\leq0.5$ with probability $0.85$ (for $\delta=0.05$), a nearly useless guarantee.

Disintegrated bounds \citep{blanchard2007occam,rivasplata2020pac,viallard2024general} directly control the risk of the drawn hypothesis: the
confidence enters through a logarithm, as in the usual bounds, instead of the multiplicative factor $1/\eta$, at the price of a
complexity term evaluated at the drawn $h$ rather than averaged over $Q$.

\bibliographystyle{plainnat}
\bibliography{../common/pacbayes}
\end{document}
